% !TEX root = ../../../main.tex
\chapter{Sets and foundations of real analysis}
\label{ch:analysis-i:sets-and-foundations}

% Shared physical dimensions for interval and number-line diagrams.
\tikzset{
  interval diagram/.style={
    x=1cm,y=6mm,font=\footnotesize,
    color=black,line width=0.5pt,line cap=round,line join=round,
    >={Stealth[length=4pt,width=3pt]}
  },
  interval endpoint/.style={
    circle,draw=black,line width=0.5pt,
    minimum size=4pt,inner sep=0pt,outer sep=0pt
  },
  interval open/.style={interval endpoint,fill=white},
  interval closed/.style={interval endpoint,fill=black},
  interval label/.style={inner sep=1pt}
}

% Reuse the physical line weights, arrowheads, and fonts of interval diagrams.
\tikzset{
  analysis diagram/.style={interval diagram,y=1cm},
  analysis guide/.style={draw=gray,densely dashed},
  analysis curve/.style={draw=black,line width=0.7pt},
  analysis label/.style={inner sep=1.5pt},
  analysis property/.style={
    draw=black,rounded corners=2pt,fill=white,
    minimum width=15mm,minimum height=8mm,inner sep=3pt
  }
}


The first three weeks move from counting sets to filling gaps in the
rational line, then from completeness to the functions built on the real
numbers. The sequence of arguments matters. Counting reveals how large
$\R$ is; completeness explains why the limits and roots we need exist.

% !TEX root = ../../../main.tex
% Sources: LEC01, IMG_9109--IMG_9111; Week 1 topic summary.
\section{Sets and countable sets}
\label{sec:analysis-i:sets-and-countable-sets}

\newthought{How do we count an infinite set?} We cannot finish listing its
elements, but we can ask whether every element has a place on a list.
Throughout these notes, $\N=\{1,2,3,\ldots\}$.

\subsection{Sets, operations, and maps}
We write $x\in A$ when $x$ belongs to $A$, and $A\subseteq B$ when every
element of $A$ belongs to $B$. Two sets are equal precisely when each is
a subset of the other. The empty set is denoted by $\varnothing$.
For sets $A$ and $B$, the basic operations are
\begin{align*}
  A\cup B&=\{x:x\in A\text{ or }x\in B\},\\
  A\cap B&=\{x:x\in A\text{ and }x\in B\},\\
  A\setminus B&=\{x:x\in A\text{ and }x\notin B\}.
\end{align*}
The complement $A^c$ is taken relative to a specified universe $U$, so
$A^c=U\setminus A$. The distributive laws and De Morgan's laws follow by
checking membership. For example,
\[
  (A\cup B)^c=A^c\cap B^c,
  \qquad (A\cap B)^c=A^c\cup B^c.
\]
For a family $(A_i)_{i\in I}$, membership in $\bigcup_{i\in I}A_i$ means
membership in at least one $A_i$; membership in $\bigcap_{i\in I}A_i$
means membership in every $A_i$.

The \emph{Cartesian product} is the set of ordered pairs
\[
  A\times B=\{(a,b):a\in A,\ b\in B\}.
\]
A map $f:A\to B$ is \emph{injective} if $f(a)=f(a')$ implies $a=a'$, and
\emph{surjective} if every $b\in B$ equals $f(a)$ for some $a\in A$.
A map with both properties is a \emph{bijection}.

\begin{definition}
  Two sets have the same \emph{cardinality}, written $|A|=|B|$, when there
  is a bijection between them. A set is \emph{countably infinite} if it
  has the same cardinality as $\N$. It is \emph{countable} if it is finite
  or countably infinite, and \emph{uncountable} otherwise.
\end{definition}
\index{countable set}\index{cardinality}
The notation $|\N|=\aleph_0$ is read ``aleph zero.'' Infinite does not
mean uncountable. For example, the even positive integers form a proper
subset of $\N$, yet $n\mapsto 2n$ is a bijection between these two sets.
\marginnote{For finite sets a proper subset has fewer elements. For infinite
sets, containment alone does not determine whether cardinality is smaller.}

\subsection{Combining countable sets}
If $A=\{a_1,a_2,\ldots\}$ and $B=\{b_1,b_2,\ldots\}$, interleave the lists
by setting
\[
  c_{2n-1}=b_n,\qquad c_{2n}=a_n.
\]
Every element of $A\cup B$ appears. If the sets overlap, discard repeated
appearances. Thus the union of two countable sets is countable, and
repeating this argument proves the same for every finite union.

What happens when there are two indices? The sum $m+n$ alone will not
distinguish $(m,n)$ from $(n,m)$. The lecture's encoding retains both pieces
of information.

\begin{proposition}
  The map $h:\N^2\to\N$ given by
  \[
    h(m,n)=m+2^{m+n}
  \]
  is injective. Consequently, the Cartesian product of two countable sets
  is countable.
\end{proposition}
\begin{proof}
  Put $d=m+n$. Since $1\leq m\leq d-1<2^d$, we have
  $2^d<h(m,n)<2^{d+1}$. These intervals are disjoint for distinct integer
  values of $d$. Hence $h(m,n)=h(j,k)$ first gives $m+n=j+k$ and then
  $m=j$, so also $n=k$. Every subset of $\N$ is countable, by listing its
  elements in increasing order. The image of $h$ therefore lists $\N^2$.
  Enumerations of two countable sets then give a list of their pairs.
\end{proof}

\begin{proposition}
  A countable union of countable sets is countable.
\end{proposition}
\begin{proof}
  Choose a list $A_m=\{a_{m,1},a_{m,2},\ldots\}$ for each nonempty set
  in the family, repeating elements if the set is finite. List the pairs
  $(m,n)$ in increasing order of $m+n$, and within each such group in
  increasing order of $m$. Every pair is reached after finitely many steps.
  Listing $a_{m,n}$ in this order reaches every element of the union.
  Discard repeated appearances to obtain the required enumeration.
\end{proof}
As usual, the argument uses a choice of enumerations for the countable
family. Figure~\ref{fig:countable-grid} shows the order of traversal.
\begin{marginfigure}[-6\baselineskip]
  \centering
  \begin{tikzpicture}[analysis diagram,x=7mm,y=7mm]
  \draw[->] (0.4,0.55) -- (4.9,0.55) node[right] {$m$};
  \draw[->] (0.55,0.4) -- (0.55,4.9) node[above] {$n$};
  \foreach \i in {1,2,3,4} {
    \node[below] at (\i,0.55) {\i};
    \node[left] at (0.55,\i) {\i};
  }
  \draw[analysis guide,->] (0.85,2.15) -- (2.15,0.85);
  \draw[analysis guide,->] (0.85,3.15) -- (3.15,0.85);
  \draw[analysis guide,->] (0.85,4.15) -- (4.15,0.85);
  \foreach \m/\n/\k in {1/1/1,1/2/2,2/1/3,1/3/4,2/2/5,3/1/6,1/4/7,2/3/8,3/2/9,4/1/10} {
    \node[fill=white,inner sep=1pt] at (\m,\n) {\k};
  }
  \node at (3.8,3.5) {$\cdots$};
\end{tikzpicture}

  \caption[Enumerating pairs along diagonals.]{Visit pairs along finite diagonals. The numerals indicate the
    order of the first ten pairs.}
  \label{fig:countable-grid}
\end{marginfigure}

\subsection{Counting the rationals}
Every positive rational has a unique reduced representation $p/q$ with
$p,q\in\N$. Thus $\Q_{>0}$ is encoded by a subset of $\N^2$ and is
countable. Adding zero and the negatives gives a countable set $\Q$.
Since $\N\subseteq\Q$, it is infinite, and therefore
\[
  |\Q|=|\N|=\aleph_0.
\]
The encoding is a map into pairs of integers; it does not assert that
rational numbers themselves are a subset of $\N$.

\subsection{Exercises on sets and enumeration}
The following exercises apply membership arguments to an arbitrary family
of sets and make the diagonal enumeration of a countable union explicit.

% !TEX root = ../../hw01.tex
% Homework 01, Exercise 1. Shared problem statement.
\begin{exercise}
  \label{ex:hw01:de-morgan}
  \solutionlink{hw01:de-morgan}
  Let $\{A_\gamma \mid \gamma \in \Gamma\}$ be a collection of sets.
  Prove De Morgan's laws
  \[
    \begin{aligned}
      \left(\bigcup_{\gamma\in\Gamma} A_\gamma\right)^c
        &= \bigcap_{\gamma\in\Gamma} A_\gamma^c,\\
      \left(\bigcap_{\gamma\in\Gamma} A_\gamma\right)^c
        &= \bigcup_{\gamma\in\Gamma} A_\gamma^c.
    \end{aligned}
  \]
\end{exercise}

% !TEX root = ../../hw01.tex
% Homework 01, Exercise 10. Shared problem statement.
\begin{exercise}
  \label{ex:hw01:countable-union}
  \solutionlink{hw01:countable-union}
  Let $A_n$ be a set such that $|A_n|=\aleph_0$ for $n=1,2,\ldots$.
  Prove $|A|=\aleph_0$ where
  \[
    A=\bigcup_{n=1}^{\infty} A_n.
  \]
\end{exercise}


% !TEX root = ../../../main.tex
% Sources: LEC01, IMG_9111--IMG_9114.
\section{Rational gaps and irrational numbers}
\label{sec:analysis-i:rational-gaps}

\newthought{Counting is only one part of the story.} The rational numbers
are closed under the usual arithmetic operations, with division by zero
excluded. Yet even a simple geometric construction leads outside $\Q$.
The right triangle in \autoref{fig:unit-triangle} has legs of length one,
so its hypotenuse has square equal to two.
\begin{marginfigure}
  \centering
  \begin{tikzpicture}[analysis diagram,x=21mm,y=21mm]
  \draw (0,0) -- node[below] {$1$} (1,0)
    -- node[right] {$1$} (1,1)
    -- node[above left] {$\sqrt{2}$} cycle;
  \draw (0.87,0) -- (0.87,0.13) -- (1,0.13);
\end{tikzpicture}

  \caption[A unit right triangle.]{A unit right triangle produces the length $\sqrt{2}$.}
  \label{fig:unit-triangle}
\end{marginfigure}

\begin{proposition}
  The number $\sqrt{2}$ is irrational.
\end{proposition}
\begin{proof}
  Suppose $\sqrt{2}=p/q$, where $p,q\in\N$ have no common factor.
  Then $p^2=2q^2$, so $p^2$ is even. An odd integer has an odd square,
  hence $p=2r$ for some integer $r$. Substitution gives $q^2=2r^2$,
  so $q$ is even as well. This contradicts the choice of $p$ and $q$.
\end{proof}

The same divisibility argument extends to other primes and higher roots.
We take up that exercise in \autoref{sec:analysis-i:roots-and-powers},
after establishing existence of the roots themselves.

\subsection{A gap in the rational line}
We can describe the location of $\sqrt{2}$ without naming a rational number
at that location. Partition $\Q$ into
\begin{align*}
  L&=\{q\in\Q\mid q\leq0\text{ or }q^2<2\},\\
  U&=\{q\in\Q\mid q>0\text{ and }q^2>2\}.
\end{align*}
Both sets are nonempty, every rational belongs to one of them, and every
element of $L$ is less than every element of $U$. There is no rational
boundary between them. Indeed, if a rational boundary $c$ satisfied $c^2<2$,
a sufficiently small positive rational increment would leave its square
below two, contradicting its being a boundary. If $c^2>2$, a sufficiently
small rational decrement gives the opposite contradiction. The case $c^2=2$
has just been ruled out. Here a boundary must be positive, since $1\in L$.
The small increments can be checked directly from
$(c+h)^2=c^2+2ch+h^2$.

This is the idea behind a \emph{Dedekind cut}. Completing the rational line
means supplying its missing boundaries. An equivalent approach supplies
limits of rational Cauchy sequences. We will formulate both versions of
completeness in \autoref{sec:analysis-i:supremum}.
\marginnote{The condition $q\leq0$ in $L$ is necessary. The sets defined only
by $q^2<2$ and $q^2>2$ would put large negative rationals on the wrong side.}

\subsection{Limits can reveal further gaps}
The lecture also introduced $e$ as an irrational number obtained from
rational factorial partial sums. This is a different source of a gap from
the geometric example above. Its construction and irrationality proof
appear in \autoref{sec:analysis-i:denseness}, where monotone convergence
and geometric-tail estimates are available. For now it motivates the
question of which limits a number system contains.

% !TEX root = ../../../main.tex
% Sources: LEC01, IMG_9114--IMG_9117.
\section{Power sets and the size of the real line}
\label{sec:analysis-i:power-sets}

The \emph{power set} $\mathcal P(X)$ is the set of all subsets of $X$.
For a finite set with $n$ elements, a subset is specified by deciding,
independently for each element, whether to include it. Thus
\[
  |\mathcal P(X)|=2^n=\sum_{k=0}^{n}\binom{n}{k}.
\]
The smallest case distinguishes the empty set from the set containing it.
% !TEX root = ../../hw01.tex
% Homework 01, Exercise 2. Shared problem statement.
\begin{exercise}
  \label{ex:hw01:power-set}
  \solutionlink{hw01:power-set}
  List all the elements in the set $\mathcal{P}(\mathcal{P}(\varnothing))$
  and verify the cardinality formula.
\end{exercise}


Cantor's theorem extends the strict increase in size to every set.

\begin{samepage}
\begin{theorem}[Cantor]
  \label{thm:analysis-i:cantor}
  There is no surjection from $X$ onto $\mathcal P(X)$.
\end{theorem}
\begin{proof}
  Let $\varphi:X\to\mathcal P(X)$ be any map, and form
  \[
    D=\{x\in X:x\notin\varphi(x)\}.
  \]
  If $D=\varphi(d)$ for some $d\in X$, then the definition gives
  $d\in D$ if and only if $d\notin D$. This is impossible, so $D$ is
  missing from the image of $\varphi$.
\end{proof}
\end{samepage}
The map $x\mapsto\{x\}$ is injective, so the theorem is expressed as
$|X|<|\mathcal P(X)|$. In particular, $\mathcal P(\N)$ is uncountable.

\subsection{Digits as subsets}
Decimal notation is a sum of place values. For example,
\[
  \pi=3+\frac{1}{10}+\frac{4}{10^2}+\frac{1}{10^3}+\cdots.
\]
More generally, a base-$b$ expansion of a number in $[0,1]$ has the form
\[
  x=\sum_{k=1}^{\infty}\frac{a_k}{b^k},
  \qquad a_k\in\{0,1,\ldots,b-1\},\quad b\geq2.
\]
Successively subdividing an interval into $b$ equal parts supplies these
digits; the remaining interval has length $b^{-n}$ after $n$ steps.
The nested interval property below justifies the limiting description.

In base two, a digit string records a subset of $\N$ through the positions
of its ones. There is a small ambiguity, such as
$0.1000\ldots=0.0111\ldots$. Choose the terminating expansion whenever
there are two choices, and represent $1$ by $0.1111\ldots$.
This gives an injection $[0,1]\to\mathcal P(\N)$.

For an injection in the other direction, send a subset $A\subseteq\N$ to
\[
  t_A=\sum_{k\in A}\frac{2}{3^k}\in[0,1].
\]
If $A$ and $B$ first differ at position $j$, the difference $2/3^j$ there
exceeds the maximum possible opposing tail
$\sum_{k>j}2/3^k=1/3^j$. Hence $t_A\ne t_B$.
The Cantor--Bernstein theorem now gives
\[
  |[0,1]|=|\mathcal P(\N)|=2^{\aleph_0}.
\]
Here we use the set-theoretic Cantor--Bernstein theorem in its usual form,
that injections each way imply a bijection; its proof is not needed for
the subsequent analysis.

The function $x\mapsto\tfrac12(1+x/(1+|x|))$ injects $\R$ into $(0,1)$.
To see injectivity, $x/(1+|x|)$ is strictly increasing on each side of zero
by cross multiplication, and has the sign of $x$. Together with the
inclusion $[0,1]\subseteq\R$, this proves the following cardinality
identity.\marginnote{The \emph{continuum hypothesis} is the additional
assertion that there is no cardinality strictly between $\aleph_0$ and
$\mathfrak c$. The identity $|\R|=2^{\aleph_0}$ does not assume that hypothesis.}
\[
  |\R|=2^{\aleph_0}=:\mathfrak c.
\]

\subsection{Exercises on the cardinality of intervals}
The next constructions compare intervals, the real line, and a square.
Here $\aleph$ denotes $|\R|=\mathfrak c$. The first construction uses the
familiar tangent function and its inverse.

% !TEX root = ../../hw01.tex
% Homework 01, Exercise 4. Shared problem statement.
\begin{exercise}
  \label{ex:hw01:interval-to-reals}
  \solutionlink{hw01:interval-to-reals}
  Obtain a bijective map from $(0,1)$ to $\mathbb{R}$.
\end{exercise}

% !TEX root = ../../hw01.tex
% Homework 01, Exercise 5. Shared problem statement.
\begin{exercise}
  \label{ex:hw01:open-to-closed-interval}
  \solutionlink{hw01:open-to-closed-interval}
  Obtain a bijective map from $(0,1)$ to $[0,1]$.
\end{exercise}

% !TEX root = ../../hw01.tex
% Homework 01, Exercise 8. Shared problem statement.
\begin{exercise}
  \label{ex:hw01:square-cardinality}
  \solutionlink{hw01:square-cardinality}
  Show that
  \[
    |(0,1)\times(0,1)|=|(0,1)|=\aleph.
  \]
  (\textit{Corrected hint.} For each number in $(0,1)$, choose the decimal
  expansion that is not eventually all $9$s. We can then write
  \[
    x\in(0,1),\quad x=0.x_1x_2\cdots,\quad
    x_i\in\{0,\ldots,9\},\quad i=1,2,\ldots.
  \]
  Then define a map from $(0,1)\times(0,1)$ to $(0,1)$ by
  \[
    (x,y)=(0.x_1x_2\cdots,0.y_1y_2\cdots)
    \mapsto (0.x_1y_1x_2y_2\ldots)=z\in(0,1).
  \]
  Show that this map is injective, then use injections in both directions.)
\end{exercise}


\subsection{Algebraic and transcendental numbers}
A real number is \emph{algebraic} if it is a root of a nonzero polynomial
with rational coefficients. It is \emph{transcendental} otherwise.
For instance, $\sqrt{2}$ is algebraic because it satisfies $x^2-2=0$.
Rational numbers are algebraic as well, since $p/q$ satisfies $qx-p=0$.
\index{algebraic number}\index{transcendental number}

\begin{proposition}
  The set $\mathcal A$ of real algebraic numbers is countably infinite.
  The set $\mathcal T=\R\setminus\mathcal A$ has cardinality $\mathfrak c$.
\end{proposition}
\begin{proof}
  For each $m\geq1$, polynomials of degree $m$ with rational coefficients
  are encoded by a subset of $\Q^{m+1}$, which is countable by repeated
  products of countable sets. Each nonzero degree-$m$ polynomial has at
  most $m$ real roots. Its factorization by $x-r$ at a root, followed by
  induction on the degree, proves this bound. Thus the set of roots of
  all such polynomials is countable. Taking the union over $m$ shows that
  $\mathcal A$ is countable. It is infinite because it contains $\Q$.

  If $\mathcal T$ were countable, then $\R=\mathcal A\cup\mathcal T$
  would be countable, contrary to Cantor's theorem. Choose distinct
  $t_1,t_2,\ldots\in\mathcal T$, and enumerate
  $\mathcal A=\{a_1,a_2,\ldots\}$ without repetition.
  Map $a_n$ to $t_{2n}$ and $t_n$ to $t_{2n-1}$, and leave every other
  transcendental number fixed. This is a bijection from $\R$ to
  $\mathcal T$, proving the cardinality assertion.
\end{proof}

Counting proves that transcendental numbers exist, without identifying a
particular one. The lecture also named $e$, $\pi$, and the Liouville number
\[
  \ell=\sum_{n=1}^{\infty}10^{-n!}.
\]
Their transcendence is stated here without proof. Including the $n=0$
term gives $\ell+1/10$, which is transcendental
as well. These stronger results belong to transcendental number theory~\cite{Waldschmidt2017}.
The earlier proof established only the irrationality of $e$.

\subsection{The Jacobian conjecture}
\label{sec:analysis-i:jacobian-conjecture}
Another lecture question concerns polynomial maps. A map
$F=(F_1,\ldots,F_d):\C^d\to\C^d$ is \emph{polynomial} when each
coordinate $F_i$ is a polynomial in $z_1,\ldots,z_d$. Its
\emph{Jacobian matrix} is the matrix of first derivatives,
\[
  \mathrm DF(z)=
  \left(\frac{\partial F_i}{\partial z_j}(z)\right)_{1\leq i,j\leq d}.
\]

\begin{conjecture}[Jacobian]
  If a polynomial map $F:\C^d\to\C^d$ satisfies
  $\det\mathrm DF(z)=c$ for every $z\in\C^d$, where $c\ne0$ is
  constant, then it has a polynomial inverse. That is, there is a
  polynomial map $G:\C^d\to\C^d$ such that
  \[
    G\circ F=F\circ G=\operatorname{id}_{\C^d}.
  \]
\end{conjecture}
This is the classical formulation~\cite{deGoursac2016}.

The determinant condition makes the linear approximation invertible at
every point. The inverse function theorem, a result from later differential
calculus, gives $F$ an inverse on a neighborhood of each point.
The conjecture asks for one inverse on all of $\C^d$, with polynomial
coordinates.

Figure~\ref{fig:jacobian-shear} illustrates a simple case,
$F(x,y)=(x+y^2,y)$, on the real plane inside $\C^2$. Here
\[
  \mathrm DF(x,y)=\begin{pmatrix}1&2y\\0&1\end{pmatrix},
  \qquad \det\mathrm DF(x,y)=1.
\]
The inverse $G(u,v)=(u-v^2,v)$ keeps the second coordinate and undoes
the horizontal shift. This verifies the conjecture for this map; the
general problem remains open~\cite{deGoursac2016}.
\begin{figure}[htbp]
  \centering
  % A finite real grid and its exact image under F(x,y)=(x+y^2,y).
\begin{tikzpicture}[analysis diagram,x=0.9cm,y=0.9cm]
  \begin{scope}
    \foreach \x in {-0.5,0,0.5} {
      \draw[gray,line width=0.5pt] (\x,-1) -- (\x,1);
    }
    \foreach \y in {-1,-0.5,0,0.5,1} {
      \draw[gray,line width=0.5pt] (-0.5,\y) -- (0.5,\y);
    }
    \draw[->] (-1,0) -- (1.15,0) node[right] {$x$};
    \draw[->] (0,-1.25) -- (0,1.35) node[above] {$y$};
    \fill (0,0) circle[radius=1.6pt];
    \node[analysis label,below left] at (0,0) {$0$};
    \fill (0,1) circle[radius=1.6pt];
    \node[analysis label,above left] at (0,1) {$(0,1)$};
    \node at (0,-1.7) {Original grid};
  \end{scope}

  \draw[->] (1.5,0.5) -- (3.65,0.5)
    node[midway,above] {$F$};
  \draw[->] (3.65,-0.5) -- (1.5,-0.5)
    node[midway,below] {$G=F^{-1}$};

  \begin{scope}[xshift=4.6cm]
    \foreach \x in {-0.5,0,0.5} {
      \draw[gray,line width=0.5pt,domain=-1:1,samples=41,variable=\t]
        plot ({\x+\t*\t},{\t});
    }
    \foreach \y in {-1,-0.5,0,0.5,1} {
      \draw[gray,line width=0.5pt]
        ({-0.5+\y*\y},\y) -- ({0.5+\y*\y},\y);
    }
    \draw[->] (-0.9,0) -- (1.95,0) node[right] {$u$};
    \draw[->] (0,-1.25) -- (0,1.35) node[above] {$v$};
    \fill (0,0) circle[radius=1.6pt];
    \node[analysis label,below left] at (0,0) {$0$};
    \fill (1,1) circle[radius=1.6pt];
    \node[analysis label,above] at (1,1) {$(1,1)$};
    \node at (0.5,-1.7) {Image under $F$};
  \end{scope}
\end{tikzpicture}

  \caption[A polynomial shear and its inverse.]{A real slice of the
    polynomial shear $F(x,y)=(x+y^2,y)$. Horizontal grid segments stay
    horizontal; vertical ones become parabolic arcs. The inverse
    subtracts $v^2$ from the first coordinate. This example has a
    global polynomial inverse.}
  \label{fig:jacobian-shear}
\end{figure}

\newthought{Motivation.} Earlier we studied injectivity
and surjectivity; a global inverse requires both. The Jacobian determinant
tests behavior near each point, so the conjecture connects those global
questions with local differential calculus. Requiring the inverse to be
polynomial adds an algebraic constraint. It illustrates why the precise
hypotheses of a theorem matter, just as existence by counting does not
identify a particular transcendental number. None of the arguments here
depends on its resolution.

\subsection{Exercises on transcendence}
The first exercise asks for the cardinality conclusion under CH, although
the argument above establishes it without that additional hypothesis.
The second applies the definition of transcendence to linear independence.

% !TEX root = ../../hw01.tex
% Homework 01, Exercise 6. Shared problem statement.
\begin{exercise}
  \label{ex:hw01:transcendental-cardinality}
  \solutionlink{hw01:transcendental-cardinality}
  Let $\aleph_0$ and $\aleph$ denote the cardinalities of $\mathbb{N}$ and
  $\mathbb{R}$ respectively. Recall that the Continuum Hypothesis (CH)
  states that there is no cardinality between $\aleph_0$ and $\aleph$.
  Use CH to prove that the set of transcendental numbers, $\mathbb{T}$,
  has the same cardinality $\aleph$ as that of $\mathbb{R}$.
\end{exercise}

% !TEX root = ../../hw01.tex
% Homework 01, Exercise 7. Shared problem statement.
\begin{exercise}
  \label{ex:hw01:linear-independence}
  \solutionlink{hw01:linear-independence}
  Let $\tau$ be a transcendental number whose existence is proven already.
  Show that for any $n\in\mathbb{N}$, the numbers
  \[
    \tau,\ldots,\tau^n
  \]
  are linearly independent over the field $\mathbb{Q}$.

  That is, if there are $r_1,\ldots,r_n\in\mathbb{Q}$ such that
  \[
    r_1\tau+\cdots+r_n\tau^n=0,
  \]
  then $r_1=\cdots=r_n=0$.

  (In the language of Linear Algebra, this result indicates that the set
  $\mathbb{R}$ viewed as a vector space over $\mathbb{Q}$ is infinitely
  dimensional.)
\end{exercise}


% !TEX root = ../../../main.tex
% Sources: LEC02, IMG_9118--IMG_9119.
\section{Convergence and iteration}
\label{sec:analysis-i:convergence-preview}

\newthought{A sequence may approach a number that its terms never reach.}
This is why a field containing all the rational numbers can still be
missing limits. We now make the meaning of ``approach'' precise.

\begin{definition}
  A sequence $(x_n)$ of real numbers \emph{converges} to $L\in\R$ if,
  for every $\varepsilon>0$, there exists $N\in\N$ such that
  \[
    n\geq N\quad\Longrightarrow\quad |x_n-L|<\varepsilon.
  \]
  We write $x_n\to L$ or $\lim_{n\to\infty}x_n=L$.
\end{definition}
The index $N$ may depend on $\varepsilon$. After that index, every term
lies in the $\varepsilon$-neighborhood $(L-\varepsilon,L+\varepsilon)$.
For convergence \emph{in $\Q$}, the limit must belong to $\Q$ as well.
The rational partial sums defining $e$ converge in $\R$, but not in $\Q$.
\marginnote{The absolute value $|x|$ is $x$ for $x\geq0$ and $-x$ otherwise.
It measures distance from zero. Its elementary inequalities are proved in
\S~\ref{sec:analysis-i:elementary-functions}.}

We will use two elementary consequences of this definition. Limits are
unique, since if $L\ne M$, the disjoint neighborhoods of radius
$|L-M|/3$ cannot both contain all sufficiently late terms. Also, if
$x_n\leq y_n$, $x_n\to x$, and $y_n\to y$, then $x\leq y$.
Otherwise neighborhoods of radius $(x-y)/3$ contradict $x_n\leq y_n$.

\subsection{Newton's idea}
Suppose we want to solve $f(x)=0$. At a current approximation $x_n$ with
$f'(x_n)\ne0$, replace the graph by its tangent line
\[
  y=f(x_n)+f'(x_n)(x-x_n).
\]
Its horizontal intercept supplies the next approximation,
\begin{equation}
  x_{n+1}=x_n-\frac{f(x_n)}{f'(x_n)}.
  \label{eq:analysis-i:newton}
\end{equation}
Figure~\ref{fig:newton-square-root} illustrates one step. This is a
motivation from calculus; a general convergence theorem for Newton's
method is not being assumed.
\begin{marginfigure}
  \centering
  \begin{tikzpicture}[analysis diagram,x=14mm,y=7mm]
  \draw[->] (-0.15,0) -- (2.8,0) node[right] {$x$};
  \draw[->] (0,-2.5) -- (0,4.2) node[above] {$y$};
  \draw[analysis curve,->,domain=0:2.4,samples=60]
    plot (\x,{\x*\x-2});
  \draw[gray,domain=1:2.5] plot (\x,{4*\x-6});
  \draw[analysis guide] (2,0) -- (2,2);
  \fill (2,2) circle[radius=1.6pt];
  \fill (1.414214,0) circle[radius=1.6pt];
  \fill (1.5,0) circle[radius=1.6pt];
  \fill (0,-2) circle[radius=1.6pt];
  \node[left] at (0,-2) {$-2$};
  \node[below left,analysis label] at (0,0) {$0$};
  \node[above left,analysis label] at (1.414214,0) {$\sqrt2$};
  \node[below,analysis label] at (1.7,-0.2) {$\tfrac32$};
  \draw[gray] (1.5,-0.1) -- (1.7,-0.35);
  \node[below,analysis label] at (2,0) {$2$};
  \node[left,analysis label] at (1.9,2.9) {$x^2-2$};
\end{tikzpicture}

  \caption[One Newton step for a square root.]{For $f(x)=x^2-2$, the tangent at $x_0=2$ meets the axis at
    $x_1=3/2$. The parabola is shown for $x\geq0$.}
  \label{fig:newton-square-root}
\end{marginfigure}

For $f(x)=x^2-a$, where $a>0$, equation~\eqref{eq:analysis-i:newton} becomes
\[
  x_{n+1}=\frac12\left(x_n+\frac{a}{x_n}\right).
\]
A positive limit $L$, if it exists, must satisfy $L^2=a$. A fixed-point
calculation alone does not establish that such a limit exists.

Here is a convergence argument using two results proved later in this
chapter. Let $r=\sqrt a>0$, whose existence follows from
Theorem~\ref{thm:analysis-i:positive-roots}, and start with $x_0>0$.
Every iterate is positive, and
\[
  x_{n+1}-r=\frac{(x_n-r)^2}{2x_n}\geq0.
\]
Thus $x_n\geq r$ for $n\geq1$, and then
$x_{n+1}-x_n=(a-x_n^2)/(2x_n)\leq0$.
Monotone convergence gives a limit $L\geq r>0$. Passing to the limit in
$2x_nx_{n+1}=x_n^2+a$ gives $L^2=a$, hence $L=r$ by uniqueness of the
positive root. If $a$ and $x_0$ are rational, every iterate is rational;
for prime $a$, the limit is irrational.

\subsection{What kind of limit can a recurrence produce?}
The lecture also considered
\[
  x_{n+1}=a+\frac{1}{x_n},\qquad a\in\N,\quad x_0>0.
\]
If this sequence converges, its limit is positive because $x_n>a$ for
$n\geq1$. The limit therefore satisfies
\[
  L^2-aL-1=0,\qquad L=\frac{a+\sqrt{a^2+4}}{2}.
\]
It is algebraic and irrational. Indeed, the rational root theorem would
force a rational root of this monic polynomial to be an integer dividing
$1$, and neither $1$ nor $-1$ is a root for $a\geq1$.
Repeated substitution also suggests a continued fraction. This example
identifies a possible limit; it does not by itself prove convergence.

Limits of rational sequences need not be algebraic. The factorial partial
sums approach $e$. Another lecture example is Euler's identity
\[
  \sum_{n=1}^{\infty}\frac{1}{n^2}=\frac{\pi^2}{6},
\]
stated here without proof~\cite{Waldschmidt2017}.
The number $\pi^2/6$ is transcendental, since a nonzero rational polynomial
vanishing there would, after substituting $X^2/6$, give a nonzero rational
polynomial vanishing at $\pi$.

\subsection{Exercise on Newton's method}
We now apply Newton's method to the fixed-point equation for the continued
fraction above. As in the square-root example, the convergence argument
uses the monotone convergence theorem proved in
\S~\ref{sec:analysis-i:completeness-equivalences}.

% !TEX root = ../../hw02.tex
% Homework 02, Exercise 6. Shared problem statement.
\begingroup
\renewcommand{\labelenumi}{(\alph{enumi})}
\begin{exercise}
  \label{ex:hw02:newton-iteration}
  \solutionlink{hw02:newton-iteration}
  In class, we discussed the iteration process
  \[
    x_{n+1}=a+\frac{1}{x_n},\quad n=0,1,2,\ldots,
    \quad a\geq 1,\quad x_0>0,
  \]
  whose limit gives us the continuous fraction
  \[
    b=a+\frac{1}{a+\frac{1}{a+\frac{1}{a+\cdots}}}.
  \]
  That is, $b$ is a fixed point of the function
  \[
    f(x)=a+\frac{1}{x},\quad x>0,
  \]
  which is realized as a root of the function $F(x)=x-f(x)$ $(x>0)$.
  \begin{enumerate}
    \item Show that the iterative algorithm
      \[
        x_{n+1}=\frac{ax_n^2+2x_n}{1+x_n^2},
        \quad n=0,1,2,\ldots,\quad x_0>0,
      \]
      arises from Newton's method that may be used to approach $b$ of our problem.
    \item Assume that $\{x_n\}$ in (a) converges.
      Verify that the limit is indeed the continuous fraction $b$ defined above.
    \item Take $a=\frac{3}{2}$. Prove $0<x_n\leq 2$ for $n=1,2,\ldots$.
    \item Prove that $\{x_n\}_{n=1,2,\ldots}$ increases for any $x_0\ne 2$
      and is constant when $x_0=2$.
    \item Show that $\{x_n\}$ converges to $2$ with any initial state $x_0$
      (thereby establishing the ``global convergence'' of the algorithm).
  \end{enumerate}
\end{exercise}
\endgroup


\subsection{Exercise on Fibonacci ratios}
The next recurrence connects an integer sequence with bounded ratios.
As in the continued-fraction example, the last part asks for a possible
limit under the assumption that the ratios converge.

% !TEX root = ../../hw03.tex
% Source: HW3 Intro to Math Analysis I (1).pdf, page 1, question 4.
% "Increasing" in part (a) is non-strict, since F_1 = F_2 = 1.
\begingroup
\renewcommand{\labelenumi}{(\alph{enumi})}
\begin{exercise}
  \label{ex:hw03:fibonacci-ratios}
  \solutionlink{hw03:fibonacci-ratios}
  The Fibonacci sequence $\{F_n\}$ is defined by
  \[
    F_0=0,\quad F_1=1,\quad F_n=F_{n-2}+F_{n-1},
    \qquad n=2,3,\ldots.
  \]
  \begin{enumerate}
    \item Show that $\{F_n\}$ is an increasing sequence.
      \emph{Here increasing means nondecreasing.}
    \item Show that $\{F_n\}$ is unbounded.
    \item Show that the Fibonacci ratios
      \[
        r_n=\frac{F_n}{F_{n-1}},\qquad n=2,3,\ldots,
      \]
      is a bounded sequence.
    \item Assume the limit of $\{r_n\}$ exists. Find the limit
      (which happens to be the golden ratio).
  \end{enumerate}
\end{exercise}
\endgroup


% !TEX root = ../../../main.tex
% Sources: LEC02, IMG_9119--IMG_9121.
\section{Fields and order}
\label{sec:analysis-i:fields-and-order}

Before asking which limits exist, we separate the algebra from the order.
A set acquires additional structure when operations are specified.
A \emph{group} is a set with an associative operation, an identity, and an
inverse for each element. It is \emph{abelian} when the operation commutes.

\begin{definition}
  A \emph{field} $K$ has two distinct elements $0$ and $1$, and operations
  of addition and multiplication, such that $(K,+)$ and $(K\setminus\{0\},\cdot)$
  are abelian groups and multiplication distributes over addition.
\end{definition}
These are structures on sets, rather than a literal containment chain of
``sets inside groups inside fields.'' The examples $\Q$ and $\R$ have their
usual arithmetic operations.

The \emph{characteristic} of a field is the least positive integer $n$ such
that $n\cdot1=0$, if there is one, and is zero otherwise. Here $n\cdot1$
means adding $1$ to itself $n$ times.

\begin{proposition}
  Every field of characteristic zero contains a unique smallest subfield,
  canonically isomorphic to $\Q$.
\end{proposition}
\begin{proof}
  Send an integer $m$ to $m\cdot1$, with the usual interpretation for
  zero and negative integers. Characteristic zero makes this map injective.
  For $m\in\Z$ and $n\in\N$, define
  \[
    \iota(m/n)=(m\cdot1)(n\cdot1)^{-1}.
  \]
  The denominator is nonzero. If $m/n=p/q$, then $mq=np$, and the same
  equality holds after mapping the integers into $K$. Multiplying by the
  inverses of the denominators shows that the definition is independent
  of the fraction chosen. The reverse argument shows injectivity.
  The usual common-denominator and product calculations show that $\iota$
  preserves addition and multiplication, so its image is a subfield.
  Any subfield of $K$ contains $1$, its integer multiples, their negatives,
  and the inverses of its nonzero elements. It therefore contains every
  $\iota(m/n)$, proving minimality and uniqueness.
\end{proof}

\subsection{Ordering the arithmetic}
A partial order is a reflexive, antisymmetric, transitive relation $\leq$.
It is a \emph{total order} if any two elements are comparable.
For example, coordinatewise order on $\Q^2$ is partial, but not total;
$(0,1)$ and $(1,0)$ are incomparable. The usual order on $\Q$ is total.

\begin{definition}
  An \emph{ordered field} is a field with a total order satisfying
  \[
    a<b\ \Longrightarrow\ a+c<b+c,
    \qquad
    a<b,\ c>0\ \Longrightarrow\ ac<bc.
  \]
\end{definition}
Every nonzero square is positive, so $1>0$. Consequently every positive
integer multiple of $1$ is positive, and an ordered field has characteristic
zero. It therefore contains the rational subfield just constructed.
The additional property distinguishing $\R$ from $\Q$ is completeness.

% !TEX root = ../../../main.tex
% Sources: LEC02, IMG_9121--IMG_9122; LEC03, IMG_9126.
\section{Supremum and completeness}
\label{sec:analysis-i:supremum}

\begin{definition}
  A number $u$ is an \emph{upper bound} of $S\subseteq\R$ if $s\leq u$
  for all $s\in S$. The \emph{supremum}, or least upper bound, is an upper
  bound $L$ such that $L\leq u$ for every upper bound $u$. We write
  $L=\sup S$. Lower bounds and the \emph{infimum} $\inf S$ are defined
  with the inequalities reversed.
\end{definition}
If a supremum belongs to $S$, it is the maximum of $S$. Membership is not
required for a supremum. For example, $(0,1)$ has supremum $1$, but no
maximum. Also, whenever the infimum exists,
$\inf S=-\sup\{-s:s\in S\}$.

We take the following completeness property as part of the description of
the real numbers.

\begin{property}[Least upper bound property (LUBP)]
  \label{thm:analysis-i:lubp}
  Every nonempty subset of
  $\R$ that is bounded above has a supremum in $\R$.
\end{property}
For a nonempty bounded-above set $S$, the assertion $L=\sup S$ is equivalent
to the following two requirements. First, every $s\in S$ satisfies $s\leq L$.
Second, for every $\varepsilon>0$, there is an $s\in S$ with
$L-\varepsilon<s$. The second requirement says that every number below
$L$ fails to be an upper bound.
\index{supremum}\index{completeness}

The rational field does not have this property. The set
\[
  S=\{q\in\Q:q\geq0,\ q^2<2\}
\]
is nonempty and bounded above, but has no rational supremum. Its supremum
in $\R$ is $\sqrt{2}$; the gap argument in
\S~\ref{sec:analysis-i:rational-gaps} excludes a rational boundary.

\subsection{Six forms of completeness}
The following properties express completeness through order, sequences,
or intervals. A sequence is \emph{bounded} if all its terms lie between
two fixed numbers. It is \emph{monotone} if it is nondecreasing or
nonincreasing. A \emph{subsequence} $(x_{n_k})$ uses strictly increasing
indices $n_1<n_2<\cdots$.

\begin{description}
  \item[LUBP] Every nonempty set bounded above has a least upper bound.
  \item[MCT] Every bounded monotone sequence converges.
  \item[NIP] Every nested sequence of nonempty closed bounded intervals
    $I_1\supseteq I_2\supseteq\cdots$ has nonempty intersection.
  \item[BWT] Every bounded sequence has a convergent subsequence
    (the Bolzano--Weierstrass theorem).
  \item[CC] Every Cauchy sequence converges. Here $(x_n)$ is \emph{Cauchy}
    when, for every $\varepsilon>0$, there is an $N$ such that
    $|x_n-x_m|<\varepsilon$ whenever $m,n\geq N$.
  \item[DC] Whenever the field is partitioned into nonempty sets $A,B$
    with $a<b$ for every $a\in A$, $b\in B$, there is a boundary $c$
    satisfying $a\leq c\leq b$ for all such $a,b$.
\end{description}
The last statement is the \emph{Dedekind cut property}. Since $c$ belongs
to one side of the partition, it is either the greatest element of $A$
or the least element of $B$. This is the precise sense in which there is
no gap at the cut.

% !TEX root = ../../../main.tex
% Sources: LEC02, IMG_9122--IMG_9125; LEC03, IMG_9126--IMG_9128.
\section{The completeness proof chain}
\label{sec:analysis-i:completeness-equivalences}

The preceding section derives the Archimedean property directly from
LUBP and proves $1/n\to0$ and $2^{-n}\to0$. Thus $\R$ has the
Archimedean hypothesis needed for the shrinking intervals in this chain.
In any Archimedean ordered field the same estimates hold.

\begin{theorem}
  In an Archimedean ordered field, the six completeness properties in
  \autoref{sec:analysis-i:supremum} are equivalent.
\end{theorem}
We follow the lecture's cycle in \autoref{fig:completeness-cycle}.
The final implication closes the argument at the beginning of week three.
The Archimedean hypothesis ensures that the bisection lengths used below
actually tend to zero.
\begin{figure*}[htbp]
  \centering
  \begin{tikzpicture}[analysis diagram]
  \node[analysis property] (lub) at (0,0) {LUBP};
  \node[analysis property] (mct) at (2.4,0) {MCT};
  \node[analysis property] (nip) at (4.8,0) {NIP};
  \node[analysis property] (bwt) at (7.2,0) {BWT};
  \node[analysis property] (cc) at (9.6,0) {CC};
  \node[analysis property] (dc) at (12,0) {DC};
  \draw[->] (lub) -- (mct);
  \draw[->] (mct) -- (nip);
  \draw[->] (nip) -- (bwt);
  \draw[->] (bwt) -- (cc);
  \draw[->] (cc) -- (dc);
  \draw[->] (dc.south) -- (12,-0.9) -- (0,-0.9) -- (lub.south);
\end{tikzpicture}

  \caption[The six completeness implications.]{The six implications. Order properties appear at the ends;
    the middle steps pass through sequences and nested intervals.}
  \label{fig:completeness-cycle}
\end{figure*}

\subsection{LUBP implies MCT}
\begin{proof}
  Let $(x_n)$ be nondecreasing and bounded above, and set
  $L=\sup\{x_n\mid n\in\N\}$. For $\varepsilon>0$, the number
  $L-\varepsilon$ is not an upper bound. There is therefore an $N$ with
  $x_N>L-\varepsilon$. For $n\geq N$, monotonicity gives
  \[
    L-\varepsilon<x_N\leq x_n\leq L<L+\varepsilon.
  \]
  Thus $x_n\to L$. Applying this argument to $(-x_n)$ proves convergence
  of every nonincreasing sequence bounded below.
\end{proof}

\subsection{MCT implies NIP}
\begin{proof}
  Write the nested intervals as $I_n=[a_n,b_n]$. Their endpoints satisfy
  \[
    a_1\leq a_n\leq a_{n+1}\leq b_{n+1}\leq b_n\leq b_1.
  \]
  By MCT, $a_n\to a$ and $b_n\to b$. Passing to limits gives $a\leq b$.
  For fixed $n$, all later endpoints lie in $[a_n,b_n]$, so
  $a_n\leq a\leq b\leq b_n$. Hence $[a,b]$ is contained in every $I_n$.
  Conversely, if $x\in I_n$ for all $n$, passage to the limit in
  $a_n\leq x\leq b_n$ gives $a\leq x\leq b$. Therefore
  \[
    \bigcap_{n=1}^{\infty}I_n=[a,b]\ne\varnothing.
  \]
\end{proof}
The intersection need not be a single point. If additionally
$b_n-a_n\to0$, then $a=b$ and the intersection is a singleton.

\subsection{NIP implies BWT}
\begin{proof}
  Let $(x_n)$ be bounded, with every term in a closed interval $I_1$ of
  positive length. Bisect $I_1$. At least one closed half contains $x_n$
  for infinitely many indices $n$. Choose such a half as $I_2$, and repeat.
  This constructs nested closed intervals $I_k$ of lengths
  $|I_1|/2^{k-1}$, each containing infinitely many indexed terms.
  By NIP, choose $L\in\bigcap_k I_k$.

  Choose $n_1$ with $x_{n_1}\in I_1$. Having chosen $n_{k-1}$, choose
  $n_k>n_{k-1}$ with $x_{n_k}\in I_k$, which is possible because infinitely
  many indices qualify. Then
  \[
    |x_{n_k}-L|\leq |I_1|\,2^{1-k}\longrightarrow0.
  \]
  Thus $(x_{n_k})$ is a convergent subsequence.
\end{proof}
\autoref{fig:nested-bisection} shows a possible sequence of choices.
\begin{marginfigure}
  \centering
  \begin{tikzpicture}[interval diagram,x=8mm]
  \foreach \a/\b/\row/\name in {0/4/0/I_1,0/2/-1/I_2,1/2/-2/I_3,1/1.5/-3/I_4} {
    \node[left] at (-0.3,\row) {$\name$};
    \draw (\a,\row) -- (\b,\row);
    \node[interval closed] at (\a,\row) {};
    \node[interval closed] at (\b,\row) {};
  }
  \node at (1.25,-4) {$\vdots$};
\end{tikzpicture}

  \caption[Nested intervals under bisection.]{Schematic bisection. Each chosen interval contains infinitely
    many terms, counted by index. The process continues beyond the rows shown.}
  \label{fig:nested-bisection}
\end{marginfigure}

\subsection{BWT implies CC}
First note that convergence always implies the Cauchy condition. If
$x_n\to L$, choose $N$ so that $|x_n-L|<\varepsilon/2$ for $n\geq N$.
For $m,n\geq N$, the triangle inequality gives
\[
  |x_n-x_m|\leq|x_n-L|+|x_m-L|<\varepsilon.
\]
The converse needs completeness.

\begin{proof}
  Let $(x_n)$ be Cauchy. Choose $N_0$ so that $|x_n-x_{N_0}|<1$ for
  every $n\geq N_0$. The tail is bounded by $|x_{N_0}|+1$, and the finitely
  many earlier terms are bounded as well. For example, the maximum of
  $|x_1|,\ldots,|x_{N_0-1}|,|x_{N_0}|+1$ bounds the whole sequence.
  BWT therefore gives a subsequence $x_{n_k}\to L$.

  Given $\varepsilon>0$, choose $N$ so that $|x_n-x_m|<\varepsilon/2$
  for $m,n\geq N$. Choose a subsequence index $n_k\geq N$ with
  $|x_{n_k}-L|<\varepsilon/2$. For every $n\geq N$,
  \[
    |x_n-L|\leq|x_n-x_{n_k}|+|x_{n_k}-L|<\varepsilon.
  \]
  Hence the full sequence converges to $L$.
\end{proof}

\subsection{CC implies DC}
\begin{proof}
  Let $A,B$ be a partition as in DC, and choose $a_1\in A$, $b_1\in B$.
  At step $n$, take the midpoint $m_n=(a_n+b_n)/2$. If $m_n\in A$,
  set $a_{n+1}=m_n$ and $b_{n+1}=b_n$; otherwise set $a_{n+1}=a_n$ and
  $b_{n+1}=m_n$. The intervals are nested, their endpoints remain on
  the prescribed sides, and
  \[
    b_n-a_n=\frac{b_1-a_1}{2^{n-1}}\longrightarrow0.
  \]
  For $m\geq n$, both $a_m$ and $a_n$ lie in $[a_n,b_n]$, so
  $|a_m-a_n|\leq b_n-a_n$. The same holds for the right endpoints.
  Consequently both sequences are Cauchy. By CC they converge, and the
  vanishing difference shows that they share a limit $c$.

  Every $a\in A$ is less than every $b_n$, so $a\leq c$ on taking limits.
  Likewise every $b\in B$ is greater than every $a_n$, so $c\leq b$.
  This is the required boundary.
\end{proof}

\subsection{DC implies LUBP}
\begin{proof}
  Let $S$ be nonempty and bounded above. Let $B$ be the set of all upper
  bounds of $S$, and let $A$ be its complement in the field. By assumption,
  $B$ is nonempty. If $s\in S$, then $s-1\in A$, so $A$ is nonempty too.
  If $a\in A$, there is an $s\in S$ with $a<s$. Thus, for every $b\in B$,
  $a<s\leq b$. The partition satisfies DC, which supplies a number $c$
  with $a\leq c\leq b$ for all $a\in A$, $b\in B$.

  We claim that $c\in B$. Otherwise some $s_0\in S$ satisfies $c<s_0$.
  Its midpoint with $c$ satisfies
  \[
    c<\frac{c+s_0}{2}<s_0.
  \]
  This midpoint cannot belong to $A$, since it exceeds $c$, and cannot
  belong to $B$, since it is smaller than the element $s_0$ of $S$.
  That contradicts the partition. Hence $c$ is an upper bound of $S$.
  Since $c\leq b$ for every upper bound $b$, it is the least upper bound.\marginnote{The
  midpoint is the useful choice. Adding a fixed amount such as $1$ to $c$
  need not leave the result below $s_0$.}
\end{proof}

\subsection{Exercise on the greatest lower bound property}
The corresponding lower-bound argument can be made directly from DC,
without passing through the least upper bound property.

% !TEX root = ../../hw03.tex
% Source: HW3 Intro to Math Analysis I (1).pdf, page 1, question 1.
\begin{exercise}
  \label{ex:hw03:dedekind-infimum}
  \solutionlink{hw03:dedekind-infimum}
  Use the existence of the Dedekind cut in $\R$ to prove the greatest lower
  bound property. That is, if $S$ is a nonempty subset of $\R$ which is
  bounded below, then there is a unique number $c\in\R$ such that
  \[
    c=\inf S.
  \]
  (Do not use LUBP.)
\end{exercise}


\subsection{Exercises on the nested interval hypotheses}
Closedness and boundedness each do work in the nested interval property.
The following examples isolate these hypotheses. The Archimedean property
from \autoref{thm:analysis-i:archimedean} supplies the indices needed
to exclude any fixed point from the intersection.

% !TEX root = ../../hw02.tex
% Homework 02, Exercise 1. Shared problem statement.
\begin{exercise}
  \label{ex:hw02:nested-bounded-intervals}
  \solutionlink{hw02:nested-bounded-intervals}
  Give an example of a sequence of nested nonempty bounded intervals
  $\{I_n\}$ in $\R$ such that $I_{n+1}\subset I_n$ for each
  $n=1,2,\ldots$ but
  \[
    \bigcap_{n=1}^{\infty} I_n=\varnothing.
  \]
  Explain what goes ``wrong''.
\end{exercise}

% !TEX root = ../../hw02.tex
% Homework 02, Exercise 2. Shared problem statement.
\begin{exercise}
  \label{ex:hw02:nested-closed-intervals}
  \solutionlink{hw02:nested-closed-intervals}
  Give an example of a sequence of nested nonempty closed intervals
  $\{I_n\}$ in $\R$ such that $I_{n+1}\subset I_n$ for each
  $n=1,2,\ldots$ but
  \[
    \bigcap_{n=1}^{\infty} I_n=\varnothing.
  \]
  Explain what goes ``wrong''.
\end{exercise}


% !TEX root = ../../../main.tex
% Sources: LEC01, IMG_9110; LEC03, IMG_9128--IMG_9130.
\section{Induction and the Archimedean property}
\label{sec:analysis-i:mathematical-induction}
\label{sec:analysis-i:archimedean}

\subsection{The natural-number tools}
Mathematical induction proves a statement $P(n)$ for every $n\in\N$ by
proving $P(1)$ and proving that $P(n)$ implies $P(n+1)$.
For example, $2^n>n$ holds at $n=1$. If it holds at $n$, then
$2^{n+1}>2n\geq n+1$, which establishes the next case. The same estimate
also follows from the binomial expansion of $(1+1)^n$.

We also use the \emph{well-ordering principle}, that every nonempty subset
of $\N$ has a least element. This is equivalent to induction. To see the
direction needed here, a subset with no least element cannot contain $1$;
if it contains none of $1,\ldots,n$, it cannot contain $n+1$ either, since
that would be its least element. Induction would make the subset empty.

\subsection{Small steps eventually pass any bound}
\begin{theorem}[Archimedean property]
  \label{thm:analysis-i:archimedean}
  For every $\varepsilon>0$ and every $M\in\R$, there is an $n\in\N$
  such that $n\varepsilon>M$.
\end{theorem}
\begin{proof}
  Suppose $S=\{n\varepsilon:n\in\N\}$ were bounded above. By LUBP,
  it would have a supremum $L$. Since $L-\varepsilon<L$, there would be
  an $m\in\N$ with $m\varepsilon>L-\varepsilon$. It would follow that
  $(m+1)\varepsilon>L$, contradicting that $L$ bounds $S$. Thus $S$ is
  unbounded above, which proves the assertion.
\end{proof}
This proof uses only LUBP, so it supplies the Archimedean hypothesis used
in the preceding completeness chain. Taking the step size to be $1$ shows
that $\N$ is unbounded in $\R$. In particular, given $\delta>0$, choose
$N>1/\delta$; then $1/n<\delta$ for all $n\geq N$, proving $1/n\to0$.

\subsection{The integer part}
\begin{proposition}
  For every $x\in\R$, there is a unique $m\in\Z$ such that
  \[
    m\leq x<m+1.
  \]
  This integer is the \emph{floor} or \emph{integer part} of $x$, denoted
  by $\lfloor x\rfloor$.
\end{proposition}
\begin{proof}
  First suppose $x\geq0$. By the Archimedean property, the set of natural
  numbers greater than $x$ is nonempty. Let $n$ be its least member.
  Then $n-1\leq x<n$, so $m=n-1$ has the required property. For a general
  real $x$, choose $k\in\N$ with $x+k\geq0$. The first case gives an
  integer $r$ with $r\leq x+k<r+1$, and $m=r-k$ works.
  Finally, two distinct integers differ by at least one, so they cannot
  both satisfy the displayed inequalities. This proves uniqueness.
\end{proof}
The lecture writes the integer part as $[x]$; we use $\lfloor x\rfloor$
to distinguish it from an interval. For example, $\lfloor\pi\rfloor=3$,
whereas $\lfloor-\pi\rfloor=-4$. The floor is not truncation toward zero.

\subsection{Exercise on bounds and convergence}
An explicit sequence connects the definitions of Cauchy convergence,
infimum, and supremum. The Archimedean property supplies the indices
needed in its estimates.

% !TEX root = ../../hw02.tex
% Homework 02, Exercise 4. Shared problem statement.
\begingroup
\renewcommand{\labelenumi}{(\alph{enumi})}
\begin{exercise}
  \label{ex:hw02:decreasing-sequence}
  \solutionlink{hw02:decreasing-sequence}
  Consider the seqeunce
  \[
    x_n=\frac{1}{n^2}+\frac{2}{n^2}+\cdots+\frac{n}{n^2},
    \quad n=1,2,\ldots.
  \]
  \begin{enumerate}
    \item Show that it is monotone decreasing and bounded below.
    \item Show that it is Cauchy.
    \item Obtain the limit of $\{x_n\}$.
    \item Let $S=\{x_1,x_2,\ldots,x_n,\ldots\}$.
      Obtain both $\inf S$ and $\sup S$.
  \end{enumerate}
\end{exercise}
\endgroup


% !TEX root = ../../../main.tex
% Sources: LEC03, IMG_9131--IMG_9133; e argument from LEC01, IMG_9111--IMG_9114.
% Geometric estimates and the positional-representation proof bridge prerequisites.
\section{Denseness, closure, and completion}
\label{sec:analysis-i:denseness}

\subsection{Upper and lower limits}
Before returning to sets, recall that $x_n\to L$ means that all sufficiently
late terms lie within every prescribed distance of $L$.
\begin{definition}
  A \emph{subsequence} of $(x_n)$ is a sequence $(x_{n_k})$ with
  $n_1<n_2<\cdots$. A real number $\ell$ is a \emph{limit point of the
  sequence} if some subsequence converges to $\ell$.
\end{definition}
We form a subsequence by keeping infinitely many terms in their original
order. The index $k$ counts the selected terms, while $n_k$ records their
positions in the original sequence. Since $n_k\geq k$, those positions
go to infinity as $k$ does. Thus a subsequence samples arbitrarily late
terms, rather than repeatedly choosing an early term.

If $x_n\to L$, every subsequence also converges to $L$. Indeed, once
all terms with $n\geq N$ lie within $\varepsilon$ of $L$, the selected
terms with $k\geq N$ do too, because $n_k\geq k\geq N$.
A convergent subsequence describes the behaviour of its selected terms.
Establishing a limit for the whole sequence requires controlling every
sufficiently late term.

A sequence may have several limit points even though the whole sequence
does not converge. For a bounded sequence, two useful quantities summarize
what remains after its first $n-1$ terms are discarded,
\[
  \alpha_n=\inf\{x_k\mid k\geq n\},\qquad
  \beta_n=\sup\{x_k\mid k\geq n\}.
\]
The sequence $(\alpha_n)$ is nondecreasing, and $(\beta_n)$ is nonincreasing.
Removing a term from a tail cannot lower its infimum or raise its supremum.
Both sequences are bounded, so their limits exist. Write
\[
  \ell_-=\liminf_{n\to\infty}x_n=\lim_{n\to\infty}\alpha_n,
  \qquad
  \ell_+=\limsup_{n\to\infty}x_n=\lim_{n\to\infty}\beta_n.
\]
These are also written $\underline{\lim}\,x_n$ and $\overline{\lim}\,x_n$.
The intervals $I_n=[\alpha_n,\beta_n]$ are nested, and $I_n$ contains
every term $x_k$ with $k\geq n$. Their endpoints approach $\ell_-$ and
$\ell_+$. Discarding finitely many initial terms changes neither limit.

\autoref{fig:tail-bounds-and-limit-points} shows this for
$x_n=(-1)^n(1+1/n)$. Choosing $n_k=2k-1$ keeps the odd terms, while
choosing $n_k=2k$ keeps the even terms. Taking $k\to\infty$ gives
\[
  x_{2k-1}=-1-\frac{1}{2k-1}\longrightarrow-1,\qquad
  x_{2k}=1+\frac{1}{2k}\longrightarrow1.
\]
The original sequence cannot converge, since its two subsequences would
then have to share the same limit. These are the only limit points,
since $|x_n|\to1$ forces
every subsequence limit $\ell$ to satisfy $|\ell|=1$.
The tail intervals shrink towards $[-1,1]$, but the numbers strictly
between $-1$ and $1$ are not limit points of this sequence.
\begin{figure*}[htbp]
  \centering
  % Exact sequence terms and tail intervals, not a schematic rescaling.
\begin{tikzpicture}[analysis diagram]
  \begin{scope}[x=4mm,y=8mm]
    \node at (6.5,2.7) {\textit{(a) Sequence terms}};
    \node at (6.5,2.2) {$x_n=(-1)^n(1+1/n)$};
    \draw[analysis guide] (0,1) -- (12.8,1);
    \draw[analysis guide] (0,-1) -- (12.8,-1);
    \draw[->] (0,-2.3) -- (0,1.9) node[above=5pt] {$x_n$};
    \draw[->] (0,0) -- (13.4,0) node[right=5pt] {$n$};
    \node[left=5pt] at (0,1) {$1$};
    \node[left=5pt] at (0,-1) {$-1$};
    \node[below left=5pt] at (0,0) {$0$};
    \foreach \n in {4,8,12} {
      \draw (\n,0.07) -- (\n,-0.07);
      \node[below=5pt] at (\n,-0.07) {$\n$};
    }
    \foreach \n in {1,...,12} {
      \fill (\n,{(-1)^\n*(1+1/\n)}) circle[radius=1.6pt];
    }
  \end{scope}

  \begin{scope}[xshift=10.7cm,x=12mm,y=8mm]
    \node at (-0.25,2.7) {\textit{(b) Bounds on successive tails}};
    \foreach \n/\lower/\upper/\row/\leftlabel/\rightlabel in {
      1/-2/1.5/1.4/{-2}/{3/2},
      3/{-4/3}/1.25/0.35/{-4/3}/{5/4},
      9/{-10/9}/1.1/-0.7/{-10/9}/{11/10}} {
      \node[anchor=east,interval label] at (-2.6,\row) {$I_{\n}$};
      \draw (\lower,\row) -- (\upper,\row);
      \node[interval closed] at (\lower,\row) {};
      \node[interval closed] at (\upper,\row) {};
      \node[above=5pt,interval label] at (\lower,\row) {$\leftlabel$};
      \node[above=5pt,interval label] at (\upper,\row) {$\rightlabel$};
    }
    \node[anchor=east,interval label] at (-2.6,-1.9) {$\bigcap_n I_n$};
    \draw[gray] (-1,-1.9) -- (1,-1.9);
    \node[interval closed] at (-1,-1.9) {};
    \node[interval closed] at (1,-1.9) {};
    \node[below=5pt,interval label] at (-1,-1.9) {$\ell_-=-1$};
    \node[below=5pt,interval label] at (1,-1.9) {$\ell_+=1$};
  \end{scope}
\end{tikzpicture}

  \caption[Tail bounds and extreme limit points.]{Terms and tail bounds
    for $x_n=(-1)^n(1+1/n)$. The first twelve terms are shown on the left.
    On the right, $I_n=[\alpha_n,\beta_n]$ contains the entire $n$th tail,
    including terms beyond those plotted. Its endpoints approach the two
    limit points $-1$ and $1$. The limiting interval bounds the limit
    points; it need not consist entirely of them.}
  \label{fig:tail-bounds-and-limit-points}
\end{figure*}

The tail infima and suprema need not be terms of the original sequence.
To show that their limits are limit points of $(x_n)$, we must select
actual terms approaching them. The following argument makes that
connection explicit for every bounded real sequence
\cite{Rodriguez2020Limsup}.
\begin{proposition}
  \label{thm:analysis-i:extreme-limit-points}
  The numbers $\ell_-$ and $\ell_+$ are respectively the smallest and
  largest limit points of a bounded real sequence $(x_n)$. Every limit
  point $\ell$ therefore satisfies $\ell_-\leq\ell\leq\ell_+$.
\end{proposition}
\begin{proof}
  Suppose $x_{n_k}\to\ell$. For each fixed $n$, all sufficiently late
  subsequence indices satisfy $n_k\geq n$. Thus
  $\alpha_n\leq x_{n_k}\leq\beta_n$. Taking limits first in $k$ and then
  in $n$ gives $\ell_-\leq\ell\leq\ell_+$.

  To reach the upper endpoint, set $n_0=0$. At step $k$, discard the
  terms through index $n_{k-1}$. The supremum of the remaining tail
  is $\beta_{n_{k-1}+1}$, so we can choose an actual term with
  $n_k>n_{k-1}$ and
  \[
    \beta_{n_{k-1}+1}-\frac1k<x_{n_k}\leq\beta_{n_{k-1}+1}.
  \]
  The indices increase, so both bounds tend to $\ell_+$. The squeeze
  theorem gives $x_{n_k}\to\ell_+$. Choosing increasing indices in the
  same way, with terms between the corresponding tail infimum and that
  infimum plus $1/k$, gives a subsequence converging to $\ell_-$.
  Hence both endpoints are limit points.
\end{proof}

For a bounded real sequence, showing that $L$ is its only limit point
is enough to establish $x_n\to L$. By
\autoref{thm:analysis-i:extreme-limit-points}, both $\ell_-$ and $\ell_+$
are limit points, so they must then equal $L$. Every original term
satisfies $\alpha_n\leq x_n\leq\beta_n$, and the squeeze theorem gives
$x_n\to L$. This is how information about subsequences controls the
whole sequence.

Equivalently, convergence occurs precisely when the two extreme limit
points agree.
If $x_n\to L$, then for every $\varepsilon>0$, all sufficiently late tails
lie in $[L-\varepsilon,L+\varepsilon]$. Their infima and suprema do too,
which proves that both limits equal $L$. The converse is the squeeze
argument just given.
This is the bounded-sequence version; more general upper and lower limits
will be studied with sequences later in the course.

The next exercise uses the elementary identities
\[
  \sin0=0,\qquad \sin\frac{2\pi}{3}=\frac{\sqrt3}{2},\qquad
  \sin\frac{4\pi}{3}=-\frac{\sqrt3}{2},
\]
together with $\sin(t+2\pi k)=\sin t$ for integers $k$.

% !TEX root = ../../hw03.tex
% Source: HW3 Intro to Math Analysis I (1).pdf, page 1, question 3.
\begin{exercise}
  \label{ex:hw03:oscillating-sequence}
  \solutionlink{hw03:oscillating-sequence}
  Consider the sequence
  \[
    a_n=\frac{1}{n}+\sin\left(\frac{2\pi n}{3}\right),
    \qquad n=1,2,\ldots.
  \]
  Find all limit points of $\{a_n\}$ and obtain
  \[
    \liminf_{n\to\infty}a_n,\quad\limsup_{n\to\infty}a_n.
  \]
\end{exercise}


\subsection{Which points can a set approach?}
\begin{definition}
  The \emph{closure} of $E\subseteq\R$ is
  \[
    \overline E=\{x\in\R\mid (x-\varepsilon,x+\varepsilon)\cap E
      \ne\varnothing\text{ for every }\varepsilon>0\}.
  \]
\end{definition}
A point of $E$ already belongs to its closure. A point $x$ is an
\emph{accumulation point} if every such neighborhood contains a point of
$E$ different from $x$. This extra condition distinguishes accumulation
points from general closure points.

\autoref{fig:closure-neighborhoods} compares these possibilities for
$E=(0,2)\cup\{3\}$. Every neighborhood of $0$ meets the interval, even
though $0\notin E$. A neighborhood of $3$ with radius less than $1$
meets $E$ only at $3$. At $5/2$, any radius less than $1/2$ misses $E$
entirely. The word \emph{every} in the definition is decisive.
\begin{figure*}[htbp]
  \centering
  % The ellipses are schematic set outlines; membership is read on the axis.
\begin{tikzpicture}[analysis diagram]
  \node at (6.8,2.1)
    {$E=(0,2)\cup\{3\},\qquad
      U_\varepsilon(x)=(x-\varepsilon,x+\varepsilon),\qquad
      0<\varepsilon_2<\varepsilon_1$};

  \foreach \shift/\testx/\smallradius/\panel in
    {0/0/0.4/1,5.3/3/0.3/2,10.6/2.5/0.3/3} {
    \begin{scope}[xshift=\shift cm]
      % The shaded lobe represents (0,2); the separate filled dot is {3}.
      \fill[gray!15] (1,0) ellipse[x radius=1cm,y radius=0.55cm];
      \begin{scope}
        \clip (1,0) ellipse[x radius=1cm,y radius=0.55cm];
        \fill[gray!40] (\testx,0)
          ellipse[x radius=\smallradius cm,y radius=0.4cm];
      \end{scope}
      \draw[gray] (1,0) ellipse[x radius=1cm,y radius=0.55cm];
      \node[analysis label] at (1.2,0.28) {$(0,2)$};

      % Two open neighborhoods with the same center and decreasing radius.
      \draw[dashed] (\testx,0)
        ellipse[x radius=0.8cm,y radius=0.85cm];
      \draw[densely dashed] (\testx,0)
        ellipse[x radius=\smallradius cm,y radius=0.4cm];
      \node[above=2pt,analysis label] at (\testx,0.85)
        {$\varepsilon_1$};
      \node[above=1pt,analysis label,fill=white] at (\testx,0.4)
        {$\varepsilon_2$};

      \draw[<->] (-0.95,0) -- (3.95,0);
      \draw[line width=1.2pt] (0,0) -- (2,0);
      \node[interval open] at (0,0) {};
      \node[interval open] at (2,0) {};
      \node[interval closed] at (3,0) {};
      \node[below=2pt,interval label] at (0,0) {$0$};
      \node[below=2pt,interval label] at (2,0) {$2$};
      \node[below=2pt,interval label] at (3,0) {$3$};

      \ifnum\panel=1
        \node at (1.5,1.5) {\textit{(a) Boundary point}, $x=0$};
        \node at (1.5,-1.18) {$0\in\overline E\setminus E$};
        \node[anchor=north,align=center,text width=4.5cm] at (1.5,-1.5)
          {Every $\varepsilon>0$ still reaches $E$.\\
           An accumulation point.};
      \else\ifnum\panel=2
        \node at (1.5,1.5) {\textit{(b) Isolated point}, $x=3$};
        \node at (1.5,-1.18) {$3\in E\subseteq\overline E$};
        \node[anchor=north,align=center,text width=4.5cm] at (1.5,-1.5)
          {Only $x$ remains when $\varepsilon<1$.\\
           Closure, but no accumulation.};
      \else
        \node at (1.5,1.5) {\textit{(c) Outside}, $x=\tfrac52$};
        \draw (2.5,-0.08) -- (2.5,0.08);
        \node[below=2pt,interval label] at (2.5,0) {$x$};
        \node at (1.5,-1.18) {$\tfrac52\notin\overline E$};
        \node[anchor=north,align=center,text width=4.5cm] at (1.5,-1.5)
          {A radius $\varepsilon<1/2$ misses $E$.\\
           One such radius excludes $x$.};
      \fi\fi
    \end{scope}
  }
\end{tikzpicture}

  \caption[Closure and shrinking neighborhoods.]{Schematic Venn views on
    a number line. The solid interval and filled point form $E$; its open
    endpoints are excluded. Dashed ellipses depict two shrinking
    $\varepsilon$-neighborhoods. The darker overlap in (a) persists as
    $\varepsilon$ decreases. Set membership is read on the axis; the
    ellipses add no points off the real line.}
  \label{fig:closure-neighborhoods}
\end{figure*}

\begin{proposition}
  A point $x$ belongs to $\overline E$ if and only if there is a sequence
  $(a_n)$ in $E$ converging to $x$.
\end{proposition}
\begin{proof}
  If $x\in\overline E$, choose $a_n\in E$ with $|a_n-x|<1/n$.
  The Archimedean property then gives $a_n\to x$. Conversely, if
  $a_n\to x$, every neighborhood of $x$ contains all sufficiently late
  $a_n$, and therefore meets $E$.
\end{proof}
For $a<b$, we have $\overline{(a,b)}=[a,b]$. Interior points already lie
in the interval, and every neighborhood of either endpoint meets it.
If $x<a$ or $x>b$, a neighborhood of radius less than its distance to the
nearest endpoint misses the interval. Thus no other point enters the closure.

\subsection{Rationals occur in every interval}
\begin{definition}
  A set $E\subseteq\R$ is \emph{dense in $\R$} if $\overline E=\R$.
  Equivalently, every nonempty open interval contains a point of $E$.
\end{definition}
For the equivalence, apply the neighborhood condition at the midpoint of
an open interval with radius half its length; in the other direction,
every neighborhood is itself an open interval.

\begin{theorem}
  \label{thm:analysis-i:rational-density}
  The rational numbers are dense in the real numbers.
\end{theorem}
\begin{proof}
  Given $x\in\R$ and $\varepsilon>0$, choose $N\in\N$ with
  $1/N<\varepsilon$. Set $m=\lfloor Nx\rfloor$. The definition of the
  floor gives $m\leq Nx<m+1$, and hence
  \[
    0\leq x-\frac{m}{N}<\frac{1}{N}<\varepsilon.
  \]
  Thus the rational number $m/N$ lies in the prescribed neighborhood
  of $x$. Every real number belongs to the closure of $\Q$.
\end{proof}
The grid in \autoref{fig:rational-grid} makes the estimate visible.
\begin{marginfigure}
  \centering
  \begin{tikzpicture}[interval diagram,x=13mm]
  \draw[<->] (-0.35,0) -- (2.7,0);
  \foreach \k in {0,1,2} {\draw (\k,-0.12) -- (\k,0.12);}
  \fill (0.65,0) circle[radius=1.6pt];
  \node[above=3pt,interval label] at (0.65,0) {$x$};
  \node[below,interval label] at (0,-0.15) {$\frac{m}{N}$};
  \node[below,interval label] at (1,-0.15) {$\frac{m+1}{N}$};
  \draw[<->] (0,-1.4) -- node[below] {$1/N$} (1,-1.4);
\end{tikzpicture}

  \caption[Approximating a real number on a rational grid.]{Schematic rational grid with spacing $1/N$. The floor chooses
    the left endpoint of the cell containing $x$.}
  \label{fig:rational-grid}
\end{marginfigure}
In particular, choosing $m_n=\lfloor nx\rfloor$ gives an explicit rational
sequence $m_n/n\to x$.

\subsection{Closure is relative to an ambient space}
Closure and completion answer related but different questions. Closure
adds points already present in a specified ambient space. Thus
$\overline{\Q}^{\,\R}=\R$, whereas $\overline{\Q}^{\,\Q}=\Q$.
A subset $E$ of an ambient space $X\subseteq\R$ is \emph{closed in $X$}
if it contains every point of its closure taken in $X$. Neighborhoods
in $X$ are intersections $(x-\varepsilon,x+\varepsilon)\cap X$.
For example, $[a,b]\cap\Q$ is closed in $\Q$. A rational point outside
it has a neighborhood missing $[a,b]$, while its real limit points
outside $\Q$ are not points of the ambient space.

A \emph{metric completion} instead embeds a metric space into a complete
one in which it is dense. For the usual distance $|x-y|$, $\R$ is a
completion of $\Q$.

More generally, $\overline E\subseteq\R$ is a completion of $E$ with its
usual distance. To check completeness, a Cauchy sequence $(y_n)$ in
$\overline E$ converges in $\R$ to some $y$. Given $\varepsilon>0$, choose
$n$ with $|y_n-y|<\varepsilon/2$, then choose $a\in E$ with
$|a-y_n|<\varepsilon/2$. The triangle inequality gives $|a-y|<\varepsilon$,
so $y\in\overline E$. Also, by the definition of closure, $E$ is dense
in $\overline E$. This explains why closing $(a,b)$ to $[a,b]$ also
completes that interval, while keeping the two notions distinct.

\subsection{Exercises on the missing limits in $\Q$}
The floor function and rational density let us build nested rational
intervals around a point missing from $\Q$. This complements the real
interval examples in \autoref{sec:analysis-i:completeness-equivalences}.

% !TEX root = ../../hw02.tex
% Homework 02, Exercise 3. Shared problem statement.
\begin{exercise}
  \label{ex:hw02:nested-rational-intervals}
  \solutionlink{hw02:nested-rational-intervals}
  Give an example of a sequence of nested nonempty bounded closed intervals
  $\{I_n\}$ in $\Q$ such that $I_{n+1}\subset I_n$ for each
  $n=1,2,\ldots$ but
  \[
    \bigcap_{n=1}^{\infty} I_n=\varnothing.
  \]
  Explain what goes ``wrong''.
\end{exercise}


\subsection{Geometric sums and factorial limits}
We need a simple estimate before the next two exercises.
Multiplying a finite geometric sum by $1-r$ and cancelling consecutive
terms gives
\[
  \sum_{k=0}^{N}r^k=\frac{1-r^{N+1}}{1-r}\qquad(r\ne1).
\]
For $0\leq r<1$, we have $r^N\to0$. If $r>0$, put $c=1/r-1>0$.
Induction gives $(1+c)^N\geq1+Nc$, and the Archimedean property shows
that $r^N\leq1/(1+Nc)\to0$. The case $r=0$ is immediate.
Taking limits in the finite identity gives
\begin{equation}
  \sum_{k=0}^{\infty}r^k=\frac1{1-r}\qquad(0\leq r<1).
  \label{eq:analysis-i:geometric-sum}
\end{equation}
Here, and in the series below, an infinite sum means the limit of its
finite partial sums. These particular sums need no general series theory.

\subsection{Completeness justifies positional notation}
We can now justify the representation facts used for cardinality in
\autoref{sec:analysis-i:power-sets}.
\begin{proof}[Proof of \autoref{thm:analysis-i:positional-representation}]
  For digits $a_k\in\{0,\ldots,b-1\}$, the partial sums
  $s_n=\sum_{k=1}^{n}a_k/b^k$ increase and satisfy
  $0\leq s_n\leq1-b^{-n}<1$ by the finite geometric identity.
  Monotone convergence supplies a limit in $[0,1]$.

  Conversely, let $0\leq x<1$, put $m_0=0$, and set
  $m_n=\lfloor b^nx\rfloor$. The floor inequalities give
  $bm_{n-1}\leq m_n\leq bm_{n-1}+b-1$.
  Thus $a_n=m_n-bm_{n-1}$ is an allowed digit, and telescoping gives
  \[
    \sum_{k=1}^{n}\frac{a_k}{b^k}=\frac{m_n}{b^n},\qquad
    0\leq x-\frac{m_n}{b^n}<b^{-n}\longrightarrow0.
  \]
  This constructs an expansion of $x$. Its digits cannot be eventually
  all $b-1$. If all digits after position $N$ were $b-1$, the geometric
  tail would give $x=(m_N+1)/b^N$, contrary to $b^Nx<m_N+1$.
  The endpoint $x=1$ is represented by taking every digit to be $b-1$.

  Finally, suppose two strings first differ at position $j$, with the
  first having the larger digit. Its advantage there is at least $b^{-j}$.
  The largest possible opposing tail is
  \[
    \sum_{k>j}\frac{b-1}{b^k}=b^{-j}.
  \]
  Equal values therefore require a digit difference of exactly one,
  followed by all zeros in the first string and all $b-1$ in the second.
  These are precisely the two forms of a terminating expansion.
  Excluding eventual strings of $b-1$ gives uniqueness in $[0,1)$.
\end{proof}

\subsection{The irrationality of $e$}
Set $s_n=\sum_{k=0}^{n}1/k!$. These partial sums increase.
For $k\geq2$, the bound $k!\geq2^{k-1}$ gives
\[
  s_n\leq2+\sum_{k=2}^{n}\frac1{2^{k-1}}<3.
\]
The monotone convergence theorem therefore supplies a real limit, denoted by
\[
  e=\sum_{k=0}^{\infty}\frac1{k!}.
\]
The alternative formula $e=\lim_{n\to\infty}(1+1/n)^n$ was also noted
in the lecture. Its equivalence with the series formula is deferred;
the series definition is enough for the present argument.

\begin{proposition}
  \label{thm:analysis-i:e-irrational}
  The number $e$ is irrational.\marginnote[-30pt]{What can we say about
    $\pi+e$ or $\pi/e$? Their irrationality is still unknown~\cite{WeissteinPi}.
    Knowing that $\pi$ and $e$ are individually transcendental does not
    settle either question.

    These questions matter because other proofs and theorems may require
    rationality or irrationality as a hypothesis. For example, irrationality
    of $\pi/e$ would imply that $a\pi+be=0$ with $a,b\in\Q$ forces
    $a=b=0$. Irrationality of $\pi+e$ would rule out $\pi+e=q$ for every
    $q\in\Q$. Such results exclude exact rational relations and supply
    hypotheses for later arguments. Until proved, any argument depending
    on these exclusions remains conditional.}
\end{proposition}
\begin{proof}
  Suppose $e=p/q$ for integers $p,q$ with $q>0$, and choose
  $n\geq\max\{q,2\}$. Both $n!e$ and $\sum_{k=0}^{n}n!/k!$ are integers.
  Their difference is
  \[
    R_n=n!\sum_{k=n+1}^{\infty}\frac1{k!}
       =\sum_{j=1}^{\infty}
         \frac1{(n+1)(n+2)\cdots(n+j)}.
  \]
  All terms are positive, and replacing each denominator factor by $n+1$
  gives, by \autoref{eq:analysis-i:geometric-sum},
  \[
    0<R_n<\sum_{j=1}^{\infty}\frac1{(n+1)^j}
          =\frac1n<1.
  \]
  The comparison is strict because the terms with $j\geq2$ are strictly
  smaller. Thus an integer lies strictly between zero and one, a
  contradiction.
\end{proof}
The factorial partial sums therefore give a bounded set of rationals
with an irrational real supremum. The next exercise asks for its bounds
and for the distinction between taking the supremum in $\R$ and in $\Q$.

% !TEX root = ../../hw02.tex
% Homework 02, Exercise 5. Shared problem statement.
\begingroup
\renewcommand{\labelenumi}{(\alph{enumi})}
\begin{exercise}
  \label{ex:hw02:factorial-sequence}
  \solutionlink{hw02:factorial-sequence}
  Consider the sequence
  \[
    x_n=\sum_{k=0}^{n}\frac{1}{k!},\quad n=1,2,\ldots,
  \]
  in $\Q$ and define $S=\{x_1,x_2,\ldots,x_n,\ldots\}$ which is a subset of $\Q$.
  \begin{enumerate}
    \item Show that $S$ is a bounded set.
    \item Find $a=\inf S$ and $b=\sup S$.
    \item Is it true that $a$ and $b$ both lie in $\Q$?
      Explain why and why not.
  \end{enumerate}
\end{exercise}
\endgroup


\subsection{Exercise on small interval covers}
The next question combines countable enumeration, rational density, and
\autoref{eq:analysis-i:geometric-sum} to construct interval covers
of small total length. Small covering length and metric incompleteness
describe different properties; completeness must be checked using Cauchy
sequences and their limits.

% !TEX root = ../../hw03.tex
% Source: HW3 Intro to Math Analysis I (1).pdf, page 2, question 5.
% The source's quoted terms "sparse" and "incompleteness" are retained.
% Small covering length alone is not a criterion for metric incompleteness.
\begingroup
\renewcommand{\labelenumi}{(\alph{enumi})}
\begin{exercise}
  \label{ex:hw03:rational-interval-cover}
  \solutionlink{hw03:rational-interval-cover}
  Let the rational numbers in $[0,1]$ be enumerated by
  $q_0,q_1,\ldots,q_n,\ldots$. Let $\varepsilon>0$ and cover this set
  $Q=\{q_n\}$ by
  \[
    U_\varepsilon=\bigcup_{n=0}^{\infty}I_n(\varepsilon),
  \]
  \[
    I_n(\varepsilon)=(q_n-\varepsilon 2^{-n},q_n+\varepsilon 2^{-n}),
    \qquad n=0,1,2,\ldots.
  \]
  That is, $Q\subset U_\varepsilon$.
  \begin{enumerate}
    \item Show that the sum of the lengths of the intervals
      $\{I_n(\varepsilon)\}$ satisfies the bound
      \[
        \sum_{n=0}^{\infty}|I_n(\varepsilon)|\leq4\varepsilon.
      \]
    \item Since $\varepsilon>0$ can be arbitrarily small, the set $Q$ can
      be covered by intervals of arbitrarily small total length indicating
      that $Q$ is a rather ``sparse'' subset of $[0,1]$. In view of this
      fact, give your explanation regarding the ``incompleteness'' of $Q$
      in $[0,1]$.
  \end{enumerate}
\end{exercise}
\endgroup


% !TEX root = ../../../main.tex
% Sources: LEC03, IMG_9133--IMG_9134; Week 3 topic summary.
\section{Polynomial functions and absolute value}
\label{sec:analysis-i:elementary-functions}

A \emph{polynomial function} has the form
\[
  p(x)=a_0+a_1x+\cdots+a_nx^n,\qquad a_0,\ldots,a_n\in\R.
\]
If $a_n\ne0$, its degree is $n$. It is defined for every $x\in\R$.
A \emph{rational function} is a quotient $p(x)/q(x)$ of polynomials,
with $q$ not the zero polynomial. Its domain excludes the zeros of $q$.
For example, the expression $(x^2-1)/(x-1)$ equals $x+1$ on its domain,
but the original quotient is undefined at $x=1$.

\subsection{Absolute value and distance}
The absolute value function, introduced before convergence in
\autoref{sec:analysis-i:convergence-preview}, is
\[
  |x|=\begin{cases}
    x,&x\geq0,\\
    -x,&x<0.
  \end{cases}
\]
It satisfies $|x|\geq0$, with equality exactly when $x=0$, and
$-|x|\leq x\leq|x|$. \autoref{fig:absolute-value} shows its graph.

\autoref{thm:analysis-i:absolute-value-inequalities} established
the product rule and triangle inequality before their use in limits.
Applying the triangle inequality to $x=(x-y)+y$, and then interchanging
$x$ and $y$, also gives the useful reverse triangle inequality
$\bigl||x|-|y|\bigr|\leq|x-y|$.

The following exercise asks for this argument and its equality case.
% !TEX root = ../../hw03.tex
% Source: HW3 Intro to Math Analysis I (1).pdf, page 2, question 7.
\begin{exercise}
  \label{ex:hw03:reverse-triangle-inequality}
  \solutionlink{hw03:reverse-triangle-inequality}
  Use the triangle inequality to prove
  \[
    \bigl||x|-|y|\bigr|\leq|x-y|,\qquad x,y\in\R.
  \]
  When will this become an equality?
\end{exercise}


\subsection{Positive and negative parts}
Define two nonnegative quantities
\[
  x^+=\max\{x,0\}=\frac{|x|+x}{2},\qquad
  x^-=\max\{-x,0\}=\frac{|x|-x}{2}.
\]
Checking separately $x\geq0$ and $x<0$ gives
\[
  x=x^+-x^-,\qquad |x|=x^++x^-,\qquad x^+x^-=0.
\]
The negative part records the \emph{magnitude} of the negative contribution,
so it is nonnegative. For a real-valued function $f$, apply these formulas
pointwise to obtain $f=f^+-f^-$ and $|f|=f^++f^-$.
\autoref{fig:positive-negative-parts} replaces the lecture's freehand
sketch with a specific example.
\begin{figure*}[htbp]
  \centering
  \begin{tikzpicture}[analysis diagram,x=21mm,y=21mm]
  \foreach \shift/\heading in {0/{$f(x)=x(1-x)$},2.2/{$f^+(x)$},4.4/{$f^-(x)$}} {
    \begin{scope}[shift={(\shift,0)}]
      \draw[->] (-0.5,0) -- (1.5,0) node[right] {$x$};
      \draw[->] (0,-0.8) -- (0,0.95) node[above] {$y$};
      \fill (0,0) circle[radius=1.6pt];
      \fill (1,0) circle[radius=1.6pt];
      \node[below left,analysis label] at (0,0) {$0$};
      \node[below left,analysis label] at (1,0) {$1$};
      \node[above] at (0.5,1.15) {\heading};
    \end{scope}
  }
  \draw[analysis curve,<->,domain=-0.45:1.45,samples=60]
    plot (\x,{\x*(1-\x)});
  \begin{scope}[shift={(2.2,0)}]
    \draw[analysis curve,<-] (-0.5,0) -- (0,0);
    \draw[analysis curve,domain=0:1,samples=40] plot (\x,{\x*(1-\x)});
    \draw[analysis curve,->] (1,0) -- (1.5,0);
  \end{scope}
  \begin{scope}[shift={(4.4,0)}]
    \draw[analysis curve,<-,domain=-0.45:0,samples=25] plot (\x,{\x*(\x-1)});
    \draw[analysis curve] (0,0) -- (1,0);
    \draw[analysis curve,->,domain=1:1.45,samples=25] plot (\x,{\x*(\x-1)});
  \end{scope}
\end{tikzpicture}

  \caption[Positive and negative parts of a function.]{Positive and negative parts of the illustrative function
    $f(x)=x(1-x)$. The part below the axis is reflected upward in $f^-$;
    in $f^+$ it is replaced by zero. All three panels use the same scale.}
  \label{fig:positive-negative-parts}
\end{figure*}

% !TEX root = ../../../main.tex
% Sources: LEC03, IMG_9135--IMG_9138; Newton argument from LEC02, IMG_9118--IMG_9119.
% Keep the section heading with its opening discussion after the wide figure.
\clearpage
\section{Roots and rational powers}
\label{sec:analysis-i:roots-and-powers}

\newthought{An inverse function needs a bijection.} In the lecture's
notation, $\R_+=[0,\infty)$ includes zero. For $n\in\N$, we will prove
that $x\mapsto x^n$ is a bijection from $\R_+$ to itself. Only after that
proof can we define $x^{1/n}$ as its inverse.

\subsection{A difference of powers}
The finite geometric identity
\[
  1+r+\cdots+r^{n-1}=\frac{1-r^n}{1-r}\qquad(r\ne1)
\]
follows by multiplying the sum by $1-r$ and cancelling consecutive terms.
For $0\leq a<b$, substitute $r=a/b$ and multiply by $(b-a)b^{n-1}$.
This gives
\begin{equation}
  b^n-a^n=(b-a)\sum_{j=0}^{n-1}b^{n-1-j}a^j.
  \label{eq:analysis-i:power-difference}
\end{equation}
The sum is positive and has $n$ terms, each at most $b^{n-1}$. Consequently,
\begin{equation}
  0<b^n-a^n\leq n(b-a)b^{n-1}.
  \label{eq:analysis-i:power-bound}
\end{equation}
In particular, the power map is strictly increasing on $[0,\infty)$,
and hence injective. The upper bound controls how much the power changes
when the input moves a short distance.

\subsection{Completeness supplies every root}
\begin{theorem}
  \label{thm:analysis-i:positive-roots}
  For every $n\in\N$ and every $y\geq0$, there is a unique $x\geq0$
  such that $x^n=y$.
\end{theorem}
\begin{marginfigure}
  \centering
  % Equal units preserve the square's geometry at its printed size.
\begin{tikzpicture}[analysis diagram]
  \fill[black!10] (0,0) rectangle (2.2,2.2);
  \draw (0,0) rectangle (2.2,2.2);
  \node[analysis label] at (1.1,1.1) {$A=s^2$};
  \node[analysis label,below=6pt] at (1.1,0) {$s$};
  \node[analysis label,left=6pt] at (0,1.1) {$s$};
  \node[analysis label,below left=5pt] at (0,0) {$(0,0)$};
  \node[analysis label,above right=5pt] at (2.2,2.2) {$(s,s)$};
\end{tikzpicture}

  \caption[A square of prescribed area.]{The square $[0,s]\times[0,s]$ in
    $\R^2$ has area $A=s^2$. Every $A\geq0$ determines a unique
    nonnegative side length.}
  \label{fig:analysis-i:root-square}
\end{marginfigure}
\begin{marginfigure}
  \centering
  % An oblique view of the solid cube; hidden edges are omitted for clarity.
\begin{tikzpicture}[analysis diagram]
  \fill[black!6] (0,1.9) -- (0.8,2.5) -- (2.7,2.5) -- (1.9,1.9) -- cycle;
  \fill[black!16] (1.9,0) -- (2.7,0.6) -- (2.7,2.5) -- (1.9,1.9) -- cycle;
  \fill[black!10] (0,0) rectangle (1.9,1.9);
  \draw (0,0) rectangle (1.9,1.9);
  \draw (0,1.9) -- (0.8,2.5) -- (2.7,2.5)
    -- (2.7,0.6) -- (1.9,0);
  \draw (1.9,1.9) -- (2.7,2.5);
  \node[analysis label] at (0.95,0.95) {$V=s^3$};
  \node[analysis label,below=6pt] at (0.95,0) {$s$};
  \node[analysis label,left=6pt] at (0,0.95) {$s$};
  \node[analysis label,below right=6pt] at (2.3,0.3) {$s$};
\end{tikzpicture}

  \caption[A cube of prescribed volume.]{The cube
    $[0,s]\times[0,s]\times[0,s]$ in $\R^3$ has volume $V=s^3$.
    Every $V\geq0$ determines a unique nonnegative side length.}
  \label{fig:analysis-i:root-cube}
\end{marginfigure}
\begin{proof}
  Uniqueness follows from strict monotonicity. If $y=0$, take $x=0$.
  Suppose henceforth that $y>0$, and put
  \[
    S=\{t\geq0\mid t^n<y\}.
  \]
  The set contains zero and also a positive element, for example
  $t=\min\{1/2,y/2\}$, since $t^n\leq t<y$.
  It is bounded above by $1+y$, because
  $(1+y)^n\geq1+y>y$ and the power map is increasing.
  LUBP therefore gives $x=\sup S>0$.

  Suppose first that $x^n<y$. Choose $h>0$ so small that
  \[
    h<1,\qquad h<\frac{y-x^n}{n(x+1)^{n-1}}.
  \]
  Using \autoref{eq:analysis-i:power-bound} with $a=x$ and $b=x+h$
  gives
  \[
    (x+h)^n-x^n\leq nh(x+h)^{n-1}
      \leq nh(x+1)^{n-1}<y-x^n.
  \]
  Hence $x+h\in S$, contradicting that $x$ is an upper bound.

  Suppose instead that $x^n>y$. Choose $h>0$ with
  \[
    h<x,\qquad h<\frac{x^n-y}{nx^{n-1}}.
  \]
  Now \autoref{eq:analysis-i:power-bound} with $a=x-h$ and $b=x$
  yields $x^n-(x-h)^n\leq nhx^{n-1}<x^n-y$, so $(x-h)^n>y$.
  Every $t\in S$ must then satisfy $t<x-h$, by strict monotonicity.
  Thus $x-h$ is an upper bound of $S$ smaller than its least upper bound,
  again a contradiction. The only remaining possibility is $x^n=y$.
\end{proof}
We may now define $y^{1/n}$ to be this unique nonnegative root.\marginnote{This
proof uses completeness and an algebraic estimate. It does not require
continuity or the intermediate value theorem, which come later in the
course. The two perturbations rule out being below and above $y$.} The inverse
is strictly increasing too. Indeed, if $u<v$ but $u^{1/n}\geq v^{1/n}$,
taking $n$th powers contradicts $u<v$.

The cases $n=2$ and $n=3$ have familiar geometric interpretations.
In $\R^2$, the square $[0,s]\times[0,s]$ has area $s^2$, as shown in
Figure~\ref{fig:analysis-i:root-square}. The theorem supplies its side
length for any prescribed area $A\geq0$. For example, a square of area
$2$ has side length $\sqrt2$.
In $\R^3$, the cube $[0,s]\times[0,s]\times[0,s]$ has volume $s^3$,
as shown in Figure~\ref{fig:analysis-i:root-cube}. A prescribed volume
$V\geq0$ likewise determines its side length, and volume $2$ gives
$s=\sqrt[3]{2}$. These pictures illustrate the root theorem rather than
replace its completeness argument.

\subsection{Exercises on roots and boundaries}
The existence of these roots lets us ask when they can be rational.
For the prime-root exercise, recall Euclid's lemma. If a prime divides
a product of integers, it divides at least one factor; repeated application
gives $p\mid a^k\Longrightarrow p\mid a$.

% !TEX root = ../../hw01.tex
% Homework 01, Exercise 3. Shared problem statement.
\begin{exercise}
  \label{ex:hw01:prime-roots}
  \solutionlink{hw01:prime-roots}
  Let $p=2,3,5,\ldots$ be a prime. Prove that $\sqrt{p}$ is irrational.
  Can this result be extended to $p^{\frac{1}{k}}$ for any integer $k\geq 2$?
\end{exercise}


% !TEX root = ../../hw01.tex
% Homework 01, Exercise 9. Shared problem statement.
\begin{exercise}
  \label{ex:hw01:irrational-roots}
  \solutionlink{hw01:irrational-roots}
  Show that if $a>0$ is irrational, so is the $k$th root of $a$,
  $a^{\frac{1}{k}}$, for any $k\in\mathbb{N}$.
\end{exercise}


The root theorem also lets us identify a Dedekind boundary by a cubic
equation. For the negative side of the cubic, $x<0$ implies $x^3<0$;
on the nonnegative side, use the strict monotonicity just established.
In the next exercise, $A$ is the upper part of the partition; the letters
are reversed from the lower-then-upper convention used in the proof chain.

% !TEX root = ../../hw03.tex
% Source: HW3 Intro to Math Analysis I (1).pdf, page 1, question 2.
% Source correction: part (a) prints B in R, but a partition member is a
% subset of R. The intended set notation is used below.
\begingroup
\renewcommand{\labelenumi}{(\alph{enumi})}
\begin{exercise}
  \label{ex:hw03:cubic-dedekind-cut}
  \solutionlink{hw03:cubic-dedekind-cut}
  Consider $A=\{a\in\R\mid a^3>5-\sqrt{2}\}$.
  \begin{enumerate}
    \item Find $B\subseteq\R$ such that $\{A,B\}$ forms a Dedekind partition.
    \item Find the Dedekind cut for this partition.
    \item Show that $A$ is bounded below and $B$ is bounded above and find
      \[
        \inf A\quad\text{and}\quad\sup B.
      \]
  \end{enumerate}
\end{exercise}
\endgroup


\subsection{A rational exponent must be well defined}
For $x\geq0$ and $m,n\in\N$, define
\[
  x^{m/n}=(x^m)^{1/n}.
\]
This agrees with $(x^{1/n})^m$, since its $n$th power is $x^m$ and it is
nonnegative. We must also check that choosing a different representation
of the fraction gives the same number.

\begin{proposition}
  If $m/n=p/q$ for positive integers $m,n,p,q$, then
  $(x^m)^{1/n}=(x^p)^{1/q}$ for every $x\geq0$.
\end{proposition}
\begin{proof}
  Let $u=(x^m)^{1/n}$ and $v=(x^p)^{1/q}$. Then
  \[
    u^{nq}=x^{mq}=x^{pn}=v^{nq},
  \]
  because $mq=pn$. Both $u$ and $v$ are nonnegative, and the $(nq)$th
  power map is injective there. Therefore $u=v$.
\end{proof}
Thus positive rational powers have a consistent meaning. At zero they
equal zero; for $x>0$, negative rational powers can subsequently be defined
by reciprocals. Real exponents require a further limiting construction.

For the next exercise, $\Q_+$ denotes the positive rational numbers,
and $\R_+=[0,\infty)$ as above. Its second part revisits the equivalent
definitions of rational powers before establishing their algebraic laws.

% !TEX root = ../../hw03.tex
% Source: HW3 Intro to Math Analysis I (1).pdf, page 2, question 6.
\begingroup
\renewcommand{\labelenumi}{(\alph{enumi})}
\begin{exercise}
  \label{ex:hw03:rational-exponent-laws}
  \solutionlink{hw03:rational-exponent-laws}
  For $x\in\R_+$, we defined the function $x^{1/n}$ for $n\in\N$.
  \begin{enumerate}
    \item For $x,y\in\R_+$, prove that $(xy)^{1/n}=x^{1/n}y^{1/n}$.
    \item For $x\in\R_+$, we defined $x^{m/n}=(x^m)^{1/n}$ for
      $m\in\N$. Show that this is the same as $(x^{1/n})^m$.
    \item For $x\in\R_+$ and $r,s\in\Q_+$, prove the laws of exponents
      \[
        (x^r)^s=x^{rs},\qquad x^{r+s}=x^rx^s.
      \]
  \end{enumerate}
\end{exercise}
\endgroup


\subsection{Exercise on odd roots of real numbers}
The preceding root construction was restricted to nonnegative inputs.
The next exercise asks for its extension to the whole real line for odd
powers.

% !TEX root = ../../hw03.tex
% Source: HW3 Intro to Math Analysis I (1).pdf, page 2, question 8.
\begin{exercise}
  \label{ex:hw03:odd-roots}
  \solutionlink{hw03:odd-roots}
  Show that for $n=3,5,7,\ldots$, the function $x^{1/n}$ is a well-defined
  monotone increasing function which is a bijection from $\R$ to itself.
  (Hint: Study $x^n$ as we did in class.)
\end{exercise}


\subsection{Returning to iteration}
We can now finish the square-root argument motivated in
\autoref{sec:analysis-i:convergence-preview}. Let $a>0$, let
$r=\sqrt a>0$ be the root supplied by
\autoref{thm:analysis-i:positive-roots}, and start with $x_0>0$.
Define $x_{n+1}=\tfrac12(x_n+a/x_n)$. Every iterate is positive, and
\[
  x_{n+1}-r=\frac{(x_n-r)^2}{2x_n}\geq0.
\]
Thus $x_n\geq r$ for $n\geq1$, and then
$x_{n+1}-x_n=(a-x_n^2)/(2x_n)\leq0$.
The monotone convergence theorem gives a limit $L\geq r>0$.
The limit laws from \autoref{thm:analysis-i:limit-laws}, applied to
$2x_nx_{n+1}=x_n^2+a$, give $L^2=a$. Root uniqueness therefore gives
$L=r$. If $a$ and $x_0$ are rational, every iterate is rational.
For prime $a$, \autoref{ex:hw01:prime-roots} shows that the limit
is irrational.

The next exercise applies the tangent slopes from
\autoref{sec:analysis-i:convergence-preview} to $F(x)=x-a-1/x$.
The root theorem, limit laws, and monotone convergence are now available
for its fixed-point and convergence arguments. Its continued fraction
denotes the positive solution of $b=a+1/b$.

% !TEX root = ../../hw02.tex
% Homework 02, Exercise 6. Shared problem statement.
\begingroup
\renewcommand{\labelenumi}{(\alph{enumi})}
\begin{exercise}
  \label{ex:hw02:newton-iteration}
  \solutionlink{hw02:newton-iteration}
  In class, we discussed the iteration process
  \[
    x_{n+1}=a+\frac{1}{x_n},\quad n=0,1,2,\ldots,
    \quad a\geq 1,\quad x_0>0,
  \]
  whose limit gives us the continuous fraction
  \[
    b=a+\frac{1}{a+\frac{1}{a+\frac{1}{a+\cdots}}}.
  \]
  That is, $b$ is a fixed point of the function
  \[
    f(x)=a+\frac{1}{x},\quad x>0,
  \]
  which is realized as a root of the function $F(x)=x-f(x)$ $(x>0)$.
  \begin{enumerate}
    \item Show that the iterative algorithm
      \[
        x_{n+1}=\frac{ax_n^2+2x_n}{1+x_n^2},
        \quad n=0,1,2,\ldots,\quad x_0>0,
      \]
      arises from Newton's method that may be used to approach $b$ of our problem.
    \item Assume that $\{x_n\}$ in (a) converges.
      Verify that the limit is indeed the continuous fraction $b$ defined above.
    \item Take $a=\frac{3}{2}$. Prove $0<x_n\leq 2$ for $n=1,2,\ldots$.
    \item Prove that $\{x_n\}_{n=1,2,\ldots}$ increases for any $x_0\ne 2$
      and is constant when $x_0=2$.
    \item Show that $\{x_n\}$ converges to $2$ with any initial state $x_0$
      (thereby establishing the ``global convergence'' of the algorithm).
  \end{enumerate}
\end{exercise}
\endgroup


\subsection{Exercise on Fibonacci ratios}
The final recurrence connects an integer sequence with bounded ratios.
Use induction for its growth and bounds, and the limit laws to identify
a possible limit. The square root appearing in that calculation now
has a rigorous construction. The last part assumes that the ratios
converge; it does not ask for a convergence proof.

% !TEX root = ../../hw03.tex
% Source: HW3 Intro to Math Analysis I (1).pdf, page 1, question 4.
% "Increasing" in part (a) is non-strict, since F_1 = F_2 = 1.
\begingroup
\renewcommand{\labelenumi}{(\alph{enumi})}
\begin{exercise}
  \label{ex:hw03:fibonacci-ratios}
  \solutionlink{hw03:fibonacci-ratios}
  The Fibonacci sequence $\{F_n\}$ is defined by
  \[
    F_0=0,\quad F_1=1,\quad F_n=F_{n-2}+F_{n-1},
    \qquad n=2,3,\ldots.
  \]
  \begin{enumerate}
    \item Show that $\{F_n\}$ is an increasing sequence.
      \emph{Here increasing means nondecreasing.}
    \item Show that $\{F_n\}$ is unbounded.
    \item Show that the Fibonacci ratios
      \[
        r_n=\frac{F_n}{F_{n-1}},\qquad n=2,3,\ldots,
      \]
      is a bounded sequence.
    \item Assume the limit of $\{r_n\}$ exists. Find the limit
      (which happens to be the golden ratio).
  \end{enumerate}
\end{exercise}
\endgroup


